\documentclass[11pt]{article}
\usepackage[margin=1in]{geometry}
\usepackage{amsmath,amssymb}
\usepackage[colorlinks=true,linkcolor=blue]{hyperref}
\newcommand{\F}{\mathbb{F}}
\newcommand{\PP}{\mathbb{P}}
\newcommand{\EE}{\mathbb{E}}
\setlength{\parindent}{0pt}
\setlength{\parskip}{5pt}

\title{Open Problems Session --- Master Index\\
\large Eighteen numbered papers (19 manuscript PDFs), 2026-06-10/28}
\author{Yitzchak Shmalo}
\date{June 15, 2026}

\begin{document}
\maketitle

\section*{Overview}

Nineteen manuscript PDFs, grouped as eighteen numbered papers, were produced in
one research session, each solving or substantially advancing a known open
problem connected to the author's work in dynamical systems, probability, random
matrix theory, and the mathematics of deep learning.  PDFs and \LaTeX{} sources are stored in the private repository
\texttt{yspennstate/math-open-problems} (one folder per paper, with generated
PDFs beside the sources where present; copies may also exist in
\texttt{Downloads}).  Five threads:

\textbf{Thread A --- non-autonomous dynamics (papers 1, 3, 4, 13).}
Source problems: the open problem in the author's May-2026 manuscript on
strict cubical horseshoes (\S14), and the genericity problem stated as open by
Balibrea--Ruck\'a (arXiv:2410.10288; Topology Appl.\ 2026).

\textbf{Thread B --- Toeplitz matrices: stability and discrepancy (papers 5,
12, 14, 16).}  Source problems: the ``Toeplitz open problem'' on asymptotic
stability of random banded Toeplitz operators (Manrique-Mir\'on's line of work),
the Matrix Spencer conjecture (Dadush--Jiang--Reis) specialized to Toeplitz and
related structured contractions, and the least-singular-value (conditioning)
question for random banded Toeplitz matrices.

\textbf{Thread C --- the uniform doubling condition (papers 2, 6--11, 15).}
Source problem: which data satisfy the uniform doubling condition of the
author's Ann.\ Math.\ Artif.\ Intell.\ (2024) paper --- left open there ---
culminating in the complete critical-scaling phase diagram and its extension
beyond the orthogonal model.

\textbf{Thread D --- implied-volatility surfaces and options mathematics
(paper 17).}  Source problem: build a theorem-level program for arbitrage-free
pricing, IV surfaces, IV percentiles, opening prices, gamma extraction, robust
pricing, and stochastic/rough-volatility models, with a concrete 100-problem
atlas and a route toward a long options-market research PDF.

\textbf{Thread E --- the extremal hard edge and phase-retrieval stability
(paper 19).}  Source problem: the author's ACHA paper on the minimum least
singular value over square row submatrices of a Gaussian matrix and the
Balan--Wang phase-retrieval stability exponential base (Randomstrasse Open
Problem 21), extended as far as possible to general random-matrix ensembles, the
three division algebras, rectangular/higher/two-sided geometries, large
deviations, a deterministic finite-alphabet injectivity obstruction, and the
hard-edge fluctuation (Gumbel) law, with the non-Gaussian upper side reduced to a
single explicit overlap-moment open problem.

\section*{The papers}

\textbf{1. \texttt{Residual\_Hilbert\_Cube\_Fusion\_v1.pdf}} (11pp).
On the Hilbert cube, depth-$N$ strict brick schemes are open and dense for
every $N$, hence residually a full hierarchy exists, with exact one-step
Poincar\'e--Miranda covering and residual entropy bursts of every order in
every tail; an obstruction theorem (the exact-horseshoe class has empty
interior) explains why the full conclusions of the \S14 problem cannot follow
from finite-depth certificates.  The remaining gap is stated precisely as an
open problem in the paper.

\textbf{2. \texttt{UDC\_High\_Probability\_v1.pdf}} (12pp).
I.i.d.\ samples with density bounded above and below on a convex body satisfy
the uniform doubling condition with probability $1-1/N$, with $N$-independent
$(\kappa,\delta,\sigma,\beta)$ and scaling factor
$\ell_N=O((\log N/N)^{1/n})$, simultaneously over all uniformly thick
truncation families.  Tools: a one-dimensional log-concave dichotomy,
Brunn--Minkowski, a cone lemma, VC relative deviations.  Two sharpness
examples included.
Lean companion: \texttt{paper2/lean/UdcAffineTransferArithmetic.lean}
machine-checks the arithmetic core of Proposition
\texttt{prop:affine-transfer}: width pullback under $t\mapsto t/a$, grid and
continuous-width transfer under $\mu_A(t)=\mu(t/a)$, and grid interpolation
with multiplier loss $\kappa\mapsto\kappa^2$.  Slab geometry, singular values,
and empirical pushforward equality remain paper-level hypotheses.

\textbf{3. \texttt{DC1\_Generic\_Convergent\_v1.pdf}} (7pp).
Resolves Balibrea--Ruck\'a's stated open problem: DC1 \emph{is} generic in
uniformly convergent surjective interval non-autonomous systems --- the DC1
class is not open, but contains the open dense class of systems whose limit
map carries a strict one-step horseshoe.  Chain: zigzag density, margin
openness, Schweizer--Sm\'{\i}tal, \v{S}tef\'ankov\'a inheritance.

\textbf{4. \texttt{Pointwise\_Baire\_Degeneracy\_v1.pdf}} (7pp).
The pointwise-metric half of the same problem: integral-type metrics make the
system spaces meager in themselves (genericity degenerates); Baire is restored
on equicontinuity strata, where the question reduces to horseshoe density ---
established conditionally, with the unconditional statement formulated as an
open problem with two identified obstructions.

\textbf{5. \texttt{Toeplitz\_Stability\_Parity\_v1.pdf}} (16pp).
The parity obstruction for random banded Toeplitz stability:
$\nu_m(\mathbb D)\equiv\mathbf 1\{G(1)G(-1)<0\}\pmod 2$ and the central limit
behavior of $(G(\pm1))$ give $\mathbb P(\text{stable})\le\frac12+o(1)$ for
every band; at odd degrees the endpoint sign has the exact finite-sample formula
\(\frac12(1-\mathbb P\{|E_q|=|O_q|\})\), and hence is exactly fair for
non-atomic laws; for Littlewood coefficients both parities are explicit
central-binomial formulas, giving
\(1/2-\sqrt{2/(\pi(m+1))}+O(m^{-1})\) at odd degree and, at even degree, exact
fairness for \(m\equiv2\pmod4\) but
\(1/2-2/(\pi m)+O(m^{-2})\) for \(m\equiv0\pmod4\); and for Gaussian
coefficients every finite degree has the exact endpoint bias
\(1/2-\pi^{-1}\arcsin((m+1)^{-1}\sum_{j=0}^m(-1)^j)\).  These endpoint
probabilities now feed a finite-degree parity envelope for every
\(\mathbb P(\nu=k)\) and every band stability event, refining the asymptotic
\(\frac12+o(1)\) cap, with explicit Gaussian/Littlewood specializations
including a strict odd-degree Littlewood stability cap
\(1/2-\sqrt{2/(\pi(m+1))}+O(m^{-1})\), and the paper now records the exact
finite-degree equality of complementary band stability probabilities;
documents a gap in a recent claimed proof of generic stability; a corrected
strong-asymmetry instability theorem; proves
\(\operatorname{Var}(\nu_m(\mathbb D))=O(m)\) for Gaussian coefficients; and
leaves the matching lower bound, sharp constant, CLT, and LLT as open problems
(see paper 5, \S6).

\textbf{6. \texttt{UDC\_Critical\_Scaling\_v1.pdf}} (7pp).
The critical-scaling dichotomy: for \emph{pure} slabs over rolling-ball bodies
the scaling factor improves to $\ell_N=O((\log N/N)^{2/(n+1)})$ (sharp
exponent), via a depth-integral slab-volume lemma (anchor dichotomy plus
Brunn--Minkowski unimodality); one thick conical truncation forces
$\ell\ge cN^{-1/n}$, so truncations shift the critical exponent from
$\frac{2}{n+1}$ to $\frac1n$.

\textbf{7. \texttt{UDC\_Log\_Necessity\_v1.pdf}} (5pp).
The logarithm is necessary: a second-moment argument over
$\asymp(N/\log N)^{(n-1)/(n+1)}$ disjoint boundary caps shows the pure-slab
condition fails with probability $1-o(1)$ below
$c(\log N/N)^{2/(n+1)}$.  With paper 6: the $k=0$ critical scaling is
$\Theta((\log N/N)^{2/(n+1)})$ exactly.

\textbf{8. \texttt{UDC\_Truncated\_Exact\_v1.pdf}} (4pp).
The truncated critical scaling is $\Theta((\log N/N)^{1/n})$ exactly: a
multi-cone second moment over $\asymp N/\log N$ orthant truncations
(overlapping bodies, disjoint apex caps).

\textbf{9. \texttt{UDC\_Edge\_Exponent\_v1.pdf}} (5pp).
A single transversal half-space truncation of the ball forces
$\ell\ge c(\log N/N)^{2/(n+2)}$ with probability $1-o(1)$ ($n\ge3$): edge
caps of mass $\asymp h^{(n+2)/2}$ along the $(n-2)$-dimensional edge.  States
the interpolation conjecture: critical exponent $\frac{2}{2n-j}$ for lowest
transversal stratum dimension $j$.

\textbf{10. \texttt{UDC\_Edge\_Exact\_v1.pdf}} (5pp).
The matching upper bound for paper 9 via a three-case section lemma for the
ball-halfspace region ($A_u(h)\ge a(n)h^{n/2}$ in every direction) and the
anchor--unimodality scheme: the edge exponent is exactly
$\Theta((\log N/N)^{2/(n+2)})$.

\textbf{11. \texttt{UDC\_Interpolation\_Law\_v1.pdf}} (7pp).
The general law, in the orthogonal model
$G_k=B^n\cap\{x_1\le0,\dots,x_k\le0\}$: for every lowest-stratum dimension
$j=n-1-k\ge1$, the critical scaling is
$\Theta((\log N/N)^{2/(2n-j)})$, both bounds with the logarithm.  Lower bound:
stratum-cap second moment on $S^j$.  Upper bound: a witness-box section lemma
whose engine is the orthogonality $\langle w,\tilde u\rangle=0$ in the
decomposition $x=y\tilde p+w$, plus an induction on $k$ trading one dimension
per truncation.  With papers 6--10 this completes the program: the exponent
ladder $\frac2{n+1},\frac2{n+2},\dots,\frac1n$ is fully proved.

\textbf{12. \texttt{Matrix\_Spencer\_Lifts\_v1.pdf}} (564pp, Thread B).
Deterministic Matrix Spencer / operator-norm discrepancy bounds for structured
self-adjoint contractions via a reusable finite-lift principle: a $K$-bounded
lift into a block-diagonal algebra
$\mathcal B=\bigoplus_\alpha\mathrm M_{r_\alpha}(\mathbb C)$ gives
$\mathrm{disc}_{\mathrm{op}}(A_1,\dots,A_n)\le
CK\sqrt{n\log(e+\mathsf E(\mathcal B)/n)}$, where
$\mathsf E(\mathcal B)=(\max_\alpha r_\alpha)\sum_\alpha r_\alpha$, as a
consequence of the block-diagonal theorem of Levy--Ramadas--Rothvoss and
Dadush--Jiang--Reis.  Instances:
scalar Toeplitz and Hankel contractions, one-level fixed-block
Toeplitz/Hankel classes, banded multilevel block-Toeplitz/Hankel classes,
matrices in a common finite-dimensional C*-algebra, and finite-group
convolution algebras with controlled irreducible dimensions.  Records the
correct entropy accounting for direct sums and entropy obstructions for exact
finite block lifts onto \(\mathrm M_D(\mathbb C)\), fixed-error faithful lifts
whose images \(\delta\)-approximate the full matrix unit ball, and
full-compressed-span finite-block reductions (entropy \(\ge D^2\) in these
full-span/full-unit-ball cases), showing that this exact or fixed-error
faithful finite-lift route alone cannot settle the square case of the full
conjecture at the class/full-unit-ball level.  The boundary section now
formalizes the tuple invariant \(\mathsf L_{K,\Lambda}\), proving discrepancy
directly in terms of finite-lift entropy and recording that structured
Toeplitz/Hankel and bounded-block classes have small \(\mathsf L\), while
full-span generic matrix tuples have \(\mathsf L=D^2\).

Current-strength upgrade folded in: Akbas--Sra's June-2026 algebraic Matrix
Spencer theorem gives the \(O(\sqrt n)\) bound for contractions in any
finite-dimensional C*-algebra with intrinsic dimension
\(\dim_{\mathbb C}\mathcal A=O(n)\).  Paper 12 now records this as an external
stage theorem, derives an AF approximation corollary for stages
\(\dim_{\mathbb C}A_k=O(n)\), and upgrades arbitrary nonabelian group
convolution in the square group-algebra regime \(|G|=O(n)\) without a
bounded-irrep hypothesis, and the perturbative version allows self-adjoint
low-Frobenius, hence rank-\(n/\log^3(en)\), noise on top of those convolution
cores.  It also folds in Akbas--Sra's low-Frobenius
perturbation theorem to get a represented-AF corollary: self-adjoint
low-Frobenius tails
\(\|L_i\|_{\mathrm F}^2\le n/\log^3(en)\) may be added to a dimension-controlled
AF-stage core at only the usual AF approximation cost; in particular the same
statement holds for operator-norm-one tails of rank at most
\(n/\log^3(en)\).  The same perturbative theorem is now packaged as a
distribution-free random ensemble result: arbitrary random algebraic cores of
intrinsic dimension \(O(n)\), including full \(D\times D\) random self-adjoint
models in the \(D^2=O(n)\) range, remain Spencer-scale with the same probability
as their structural event after adding arbitrary dependent low-Frobenius,
low-rank, spiked, compressed, or coordinate-supported self-adjoint tails.  The
random section now starts with a distribution-free all-plus gauge for arbitrary
self-adjoint or rectangular random matrices: after pathwise scaling by
\[
  \max\left\{\max_i\|X_i\|_{\rm op},
  n^{-1/2}\left\|\sum_iX_i\right\|_{\rm op}\right\},
\]
the all-plus signing gives discrepancy at most \(\sqrt n\), and active-menu
deterministic scales only have to dominate this gauge eventually.  A new
moment-schedule corollary makes this deterministic-scale check explicit:
\(\sum_{m,h}\mathbb E(G_{m,h}/s_{m,h})^p<\infty\) already gives eventual
almost-sure contraction normalization and the same all-plus certificate, with
rowwise quantile budgets as a direct specialization; the closed-form schedule
version gives all-integer and sparse-dimension thresholds when the active menu
size is polynomial in the ambient dimension.  A tail-profile schedule version
accepts direct sub-Weibull or user-supplied gauge tails and gives logarithmic
deterministic normalizations, again without independence.  A finite-confidence
version gives one-row or finite-library certificates with explicit
\(\delta\)-dependent scales, and the conditional version lets those scales be
chosen from side information through conditional moment or tail envelopes.
Random active finite-confidence menus also pay only their selected conditional
error budget rather than a deterministic ambient library size.  The same
conditional budgets now have a Borel--Cantelli schedule form, giving adaptive
almost-sure normalizations from summable conditional moment or tail budgets,
with the same selected-budget rule for random active submenus.  It also
records a full-matrix average-case route: centered, possibly dependent,
Gaussian-operator-tail ensembles stable under scalar \(\ell_2\) sums, including
GOE/GUE and centered sub-Gaussian Wigner-type models, satisfy Matrix Spencer in
every ambient dimension after deterministic scaling by
\(K(\sqrt D+\sqrt{\log(en)})\), since the all-plus signing already has norm at
most \(\sqrt n\) with high probability; random active operator-tail libraries
pay only their selected expected tail count, with the same selected-cost
principle for rectangular/non-self-adjoint operator-tail libraries.  A sign-symmetry transfer now shows
that for coordinatewise sign-invariant random tuples any all-plus
\(O(\sqrt n)\) tail is automatically a discrepancy tail for the same
realization-normalized ensemble; its Borel--Cantelli active-menu form gives
simultaneous eventual signings for finite or countably exhausted
sign-symmetric model libraries whenever the all-plus failure series is
summable.  Independent and conditionally independent bounded sign-symmetric
arrays get a concrete matrix-Bernstein active-menu criterion from the all-plus
variance proxy, with summable environmental envelope failures, all-integer and
sparse dimension schedules under logarithmic-square sample growth, and
closed-form discrepancy-margin constants; symmetric scalar-multiplier
environments inherit these explicit constants from the conditional
second-moment profile \(\|\sum_iq_iB_i^2\|_{\mathrm{op}}\), and unbounded
symmetric sub-Gaussian multipliers are covered after cap normalization with the
matrix-Laplace margins \(K>\sqrt{2\beta(\Gamma+1)}\) and
\(K>\sqrt{2\beta\Gamma}\).  Arbitrary symmetric multiplier tails are covered by
a tail-capped companion whenever cap failures are summable; sub-Weibull tails
pay \(2\sum_i e^{-(L/a_i)^\alpha}\) and retain the bounded-Bernstein margins
\(\Theta(\Gamma+1,a,\beta)\) and \(\Theta(\Gamma,a,\beta)\).  Common-exponent
sub-Weibull operator tails with
\(\mathbb P\{\|X_i\|_{\mathrm{op}}>t\mid\mathcal P_i\}\le
A\exp\{-(t/s)^\alpha\}\) now have the direct dominated schedule
\(r=s(\log(e nA/\varepsilon))^{1/\alpha}\) and
\(V=C_\alpha nAs^2\).  Polynomial
moment proxies of any order \(p>0\) are handled explicitly by writing
\(r=p\wedge2\), paying cap cost \(\sum_i(a_i/L)^p\), and controlling the
truncated variance by \(L^{-r}\|\sum_i a_i^rB_i^2\|_{\mathrm{op}}\); the same
criteria allow environment-measurable adaptive caps and moment proxies before
the multipliers are sampled.  Common finite-moment operator tails also have a
dominated-increment schedule: if
\(\mathbb E(\|X_i\|_{\mathrm{op}}^p\mid\mathcal P_i)\le As^p\), then
\(S\le nAs^p\) and, for \(p\ge2\), \(W\le nA^{2/p}s^2\), yielding explicit
\(n\)- and \(D\)-power normalizations from the common-moment optimizer.  A mixed
active-library corollary lets active streams in the same row be tagged either
sub-Weibull or finite-moment, with per-stream exponents and the same cap-budget
and matrix-tail summability, so countable mixed dominated-tail libraries are
exhausted by one active prefix.  A signed-core companion lets any active-prefix
core library with one common signing absorb these mixed dominated perturbations
after per-stream scaling, with the same all-integer and sparse thresholds.  A
raw noncentered version keeps matrices \(W=B+Y\) together, normalizes by
\(\max\{1,\lambda\}\), and therefore covers structured conditional-mean cores
plus mixed dominated centered fluctuations directly; a scaled version replaces
the unit structured-mean bound by an environment-measurable size \(\rho\) and
normalizes by \(\max\{\rho,\lambda\}\), while a two-scale version decouples the
row size \(\rho\) from the signed-core size \(\sigma\) and uses
\(\max\{\rho,\sigma,\lambda\}\); finite-lift and algebraic structured cores now
plug into this raw two-scale form directly.  A
finite-lift specialization now packages that directly for active-prefix
finite-lift core menus, replacing the full Gaussian operator-tail scale by
per-stream mixed dominated \(q_{m,h}\) normalizations.  An algebraic-core menu
specialization uses Akbas--Sra on the direct-sum active algebra, so active menus
with total intrinsic C*-dimension \(O(n)\) absorb the same mixed dominated
perturbations.  A catalogue corollary records the named cases explicitly:
Toeplitz/Hankel, banded multilevel, finite C*-algebra, and finite group
convolution core menus are covered under the corresponding active-prefix
entropy or intrinsic-dimension budgets.  A concrete size-budget version makes
the tests explicit through \(\sum_h D_{m,h}\),
\(\sum_h\prod_j(2w_{m,h,j}+1)\),
\(\sum_h\dim_{\mathbb C}\mathcal A_{m,h}\), or \(\sum_h |G_{m,h}|\)
bounded by \(O(n)\).  A random-orbit closure version preserves the same
mixed-tail conclusion under random unitary or permutation conjugacy, corners,
fixed-size tensor amplification, block direct sums, and bounded self-adjoint
images with the corresponding summed charge.  A selected-submenu version lets
the active structured library be chosen randomly or adaptively from the row
environment when its cardinality and realized total charge are summably
controlled.  A Bernstein-selected version checks those selection failures from
independent candidate indicators using deterministic probability and charge
envelopes; a Laplace-dominated version covers negatively dependent selections
through product-mgf bounds; a random-envelope Laplace version permits
environment-measurable probability envelopes once aggregate cardinality, charge,
and quadratic-charge profiles are summably certified; a random-envelope
Poisson-count version covers proposal-count and occupancy catalogues through
compound-Poisson cardinality and charge profiles; a Cox-Poisson
occupied-support refinement charges labels by \(1-e^{-\Lambda}\), so repeated
proposals do not increase the occupied-support budget; a product-categorical
version handles randomized one-template-per-slot library choices with
deterministic cardinality and a scalar Bernstein charge tail; an
adaptive-categorical version allows sequential
one-template-per-slot choices under predictable probability envelopes and a
Freedman charge tail; a random-envelope adaptive-categorical version allows
state-dependent predictable envelopes once their aggregate charge drift and
quadratic charge are summably certified; a Freedman-selected version removes
conditional independence for adaptive sequential indicators with predictable
probability envelopes; a random-envelope Freedman-selected version makes those
predictable envelopes random once aggregate cardinality, charge, and
quadratic-charge profiles are summably certified; and a Serfling-selected
version handles fixed-size and stratified sampling without replacement from
finite-population charge averages.  A hybrid composition corollary combines
these selection mechanisms componentwise with no cross-mechanism independence,
provided the component failure, cardinality, and charge budgets are summably
certified; a random-budget hybrid version allows the component budgets
themselves to be environment-random after summable aggregate certification.
A hierarchical random-budget version flattens recursive random trees, branching
catalogues, adaptive search frontiers, and random forests of structured cores
to selected terminal antichains with summably certified frontier size and
charge.
The two
random routes combine:
any already-signed random core class, including dimension-controlled algebraic,
finite-lift, Toeplitz, and Hankel cores, remains Spencer-scale after adding
full random operator-tail noise whose scalar linear combinations obey the same
Gaussian tail conditionally on the core, in arbitrary ambient dimension after
the same natural spectral scaling of the noise.  For active-prefix finite-lift
core libraries this is simultaneous after the union cost
\(J_m(n_m+1)e^{-\tau_m^2}\) and noise scaling \(\sqrt{D_m}+\tau_m\);
the explicit schedules are arbitrary-growth
\(\tau_m^2\gtrsim\log(J_m(n_m+1))+\log m\), all-integer
\(\tau_D^2=R\log(eD)\) with \(R>\Gamma+1\), and sparse \(R>\Gamma\)
when \(J_m(n_m+1)\le(eD_m)^\Gamma\).
A rank-sensitive matrix-series specialization improves this perturbation scale
to \(q_{m,h}(\sqrt{\log(e\kappa_{m,h})}+\tau_m)\) for conditionally Gaussian,
Rademacher, or scalar sub-Gaussian coefficient series with predictable
self-adjoint directions, realized coefficient square-function bounds
\(\max_i\|\sum_rG_{i,r}^2\|\le q_{m,h}^2\),
\(\|\sum_{i,r}G_{i,r}^2\|\le n_mq_{m,h}^2\), and rank envelope
\(\kappa_{m,h}\); the full-rank fallback is \(\kappa_{m,h}=D_m\).
A common-projection specialization makes this directly checkable for low-rank,
spiked, block-localized, and random-subspace matrix-series noise: once all
coefficient directions live in a rank-\(\kappa_{m,h}\) projection, only the two
square-function bounds remain.
An atomic norm-budget specialization reduces those two checks to scalar
amplitude sums, covering rank-one spike fields, self-adjoint edge atoms, and
finite frame/dictionary matrix-series noise inside the same low-rank support.
A block-load version lets atoms touch several orthogonal blocks and pays only
the maximum incident block energy, while the logarithmic factor sees the total
rank of the active block family.
Bounded weighted block-designs, including regular, matching, permutation, and
balanced without-replacement designs, need no load-concentration payment once
their weighted block degrees are deterministic or fail only summably.
A bounded-overlap block-hypergraph specialization gives the direct
\(q/\sqrt d\) rule for block designs, regular block hypergraphs, and balanced
block dictionaries with overlap degree \(d\).
A heterogeneous degree-profile version uses local row-block and total-block
degrees, avoiding worst-block normalization for irregular block hypergraphs.
A row-wise block-permutation specialization gives the concrete \(q/\sqrt d\)
amplitude normalization for \(d\)-round permutation, matching, and balanced
without-replacement assignments.
Conditional scalar Bernstein verifies these block-load budgets for independent
random block choices, random edge-atom supports, and clustered frame
dictionaries, and a weighted version allows random atom amplitudes,
compound-thinned atom fields, and random weighted dictionaries.
A uniform block-support specialization gives explicit budgets for atoms choosing
fixed-size random block sets, Bernoulli block thinnings, or random block
hyperedges.
A profile schedule form turns the resulting Bernstein exponents into
all-integer and sparse Borel--Cantelli thresholds.
Its homogeneous form reduces those exponents to the expected number of active
atom incidences per block.
A predictable block-support profile form removes independence for sequential
support exposure by using one-step inclusion probabilities and predictable
amplitude envelopes.
Its profile schedule has the same active block-count logarithm as the weighted
Freedman schedule, with the predictable-support exponents in place of the
weighted incident-load exponents.
An inhomogeneous incidence-mass specialization replaces the raw exponents by
row-block and total-block predictable probability masses, retaining the same
\(8/3\) threshold constants for biased adaptive block samplers.
Its homogeneous predictable uniform schedule has the same
\(8(\Gamma+1)/3\) all-integer and \(8\Gamma/3\) sparse incidence thresholds as
the independent homogeneous block-support schedule.
A predictable weighted version replaces independence by scalar Freedman budgets,
covering adaptive atom sampling, online random weighted dictionaries, and
without-replacement block exposure.
Both weighted block-load variants now have profile schedules once their
Bernstein or Freedman exponents dominate the active block count.
A bounded weighted incidence specialization needs only the normalized
row-block and total-block mean loads \(\rho/\omega\), with the same \(8/3\)
threshold constants.
A quadratic-form-to-operator-tail criterion verifies this input from conditional
scalar sub-Gaussian bounds on \(\langle u,\sum_i a_iY_i u\rangle\), so
conditionally Gaussian, correlated Gaussian, and Wigner-type active-prefix noise
libraries enter with scale \(q_{m,h}(\sqrt{D_m}+\tau_m)\).
For conditionally independent upper-triangular sub-Gaussian entries, the proxy is
the entry square-function envelope \(Q_{m,h}\), and uniform entry scale
\(\sigma_{m,h}\) gives the usual
\(\sigma_{m,h}(\sqrt{D_m}+\tau_m)\) Wigner/GOE/GUE normalization.
A row-variance envelope gives sparse and banded entry-profile libraries the
explicit scale
\((\alpha^2+2\Delta\beta^2)^{1/2}(\sqrt{D_m}+\tau_m)\).
At normalized graph scale \(\beta\lesssim\Delta^{-1/2}\), this keeps the same
all-integer and sparse-dimension active-prefix schedules without degree loss.
Independent random graph masks are now included after a scalar Bernstein
max-degree exception; inhomogeneous Erdos--Renyi supports use expected-degree
normalization once the expected degree dominates the active-prefix union
logarithm, while random regular, matching, and other dependent bounded-degree
masks need no support exception.  A weighted random-mask variant replaces degree
by a scalar-Bernstein bound for the realized row-variance profile
\(2\sum_q\eta_{pq}c_{pq}^2\), covering highly inhomogeneous variance-profile
entry masks; its schedule corollary gives arbitrary-growth, all-integer, and
sparse-dimension thresholds when the support Bernstein exponent dominates the
active-prefix union logarithm.  A bounded-incidence specialization reduces this
to the normalized mean row variance \(\rho/\omega\), retaining the same \(8/3\)
threshold constants.  A predictable-profile extension allows the
inclusion probabilities, proxy caps, and Bernstein budgets themselves to be
chosen from the core environment, paying only a summable envelope-exception
probability; adaptive and randomly sampled variance-profile masks therefore use
the same schedules.  A dependent weighted-budget mask corollary removes support
independence altogether once the realized weighted row budgets are controlled,
so random regular, matching, without-replacement, determinantal, and other
correlated masks enter at the same finite-lift scale.  Its bounded
weighted-degree specialization gives the usual \(d^{-1/2}\) edge normalization
for arbitrary dependent weighted-degree-bounded masks with no support
exception.
Full random matrix ensembles with scalar-combination Gaussian operator tails
also have a direct active-menu all-plus certificate: after normalization by
\(K_{m,h}(\sqrt{D_m}+\tau_m)\), the all-plus signing works simultaneously for
finite or countable libraries whenever \(J_m(n_m+1)e^{-\tau_m^2}\) is
summable, with no sign-symmetry assumption.  A quadratic-form net criterion
makes the self-adjoint hypothesis checkable from scalar tails for
\(\langle u,\sum_i a_iY_i u\rangle\); Hermitian sub-Gaussian coefficient
series are controlled by their atom square functions.  Independent
upper-triangular Hermitian entries have an explicit quadratic-form square
function, bounded by \(\sqrt2\,\sigma\) in the homogeneous case and by row
variance budgets for sparse or banded profiles.  A Bernstein quadratic-form
variant covers centered rank-one covariance and averaged Wishart blocks at scale
\(\kappa^2((\sqrt D+\tau)/\sqrt r+(\sqrt D+\tau)^2/r)\), and Hermitian
independent entries with Bernstein variance/jump proxies at scale
\(Q(\sqrt D+\tau)+B(\sqrt D+\tau)^2\).  The same active all-plus
certificate covers rectangular and non-self-adjoint full operator-tail matrices
with \(D_m^{\rm rec}=d_{m,1}+d_{m,2}\).  A bilinear-form net lemma makes this
checkable from scalar sub-Gaussian tails for
\(\langle u,\sum_i a_iX_i v\rangle\), with selected-cost random-active forms
for environment-selected bilinear libraries, at scale
\(q_{m,h}(\sqrt{D_m^{\rm rec}}+\tau_m)\).  Haar unitary/orthogonal active menus
have \(q=O(D^{-1/2})\), so they need only a constant normalization when
\(\tau=O(\sqrt D)\).  Independent-entry rectangular
sub-Gaussian arrays have \(q_{m,h}\) bounded by the worst weighted bilinear
entry proxy, in particular by \(\max_{i,r,s}\sigma_{m,h,i,r,s}\).  More general
rectangular sub-Gaussian coefficient series are controlled by the worst
fixed-bilinear atom energy, bounded by the left/right atom variance operators.
A Bernstein bilinear-form variant gives the same rectangular all-plus
certificate at scale
\(q(\sqrt{D_m^{\rm rec}}+\tau_m)+b(\sqrt{D_m^{\rm rec}}+\tau_m)^2\), with the
same selected-cost random-active form for environment-selected Bernstein
bilinear libraries, covering bounded, Poissonized, exponential-entry, and
gamma-entry rectangular libraries and selected independent-entry rectangular
arrays through explicit variance and jump proxies; conditionally independent
entrywise Bernstein mgf assumptions now imply those proxies directly, with a
one-row finite-confidence version for finite selected libraries,
and bounded-entry finite-confidence certificates require only conditional
second-moment proxies plus uniform entry caps,
and balanced Haar rank-one dyads after only constant amplitude normalization.
A bounded matrix-Bernstein/Freedman all-plus route removes sign symmetry for
independent or adapted centered bounded self-adjoint and rectangular blocks
under summable ordinary or predictable all-plus variance budgets, covering
centered random projections, bounded sparse increments, finite-confidence
independent and adapted selected libraries, online feature blocks,
adaptive measurement increments, and martingale-difference matrix streams.
A quantum-effect finite-confidence specialization covers centered random POVM
effects and density matrices, with variance controlled by
\(\|\sum_i\mathbb E(E_i\mid\mathcal F)\|_{\rm op}\).
An explicit schedule corollary turns the half-budget profile
\(v\le\beta s^2n/\log(e\Delta)\) into active-menu, all-integer, and sparse
thresholds through
\(\gamma(K,\beta,B)=K^2/(2(\beta+KB/3))\).
A proxy-scale variant separates the typical variance scale \(q\) from the
largest jump \(b\): if \(b\le Rq\) and
\(v\le\beta q^2n/\log(e\Delta)\), the same active schedules hold after
normalization by \(\max\{R,K\}q\), with
\(\gamma_{\rm pr}=K^2/(2(\beta+KRB/3))\).
A tail-capped centered Bernstein/Freedman variant now covers unbounded
independent or adapted self-adjoint and rectangular streams: after cap failures
are summable, the capped centered ordinary or predictable variance obeys the
same half-budget profile, and the truncated mean drift is at most
\(\eta q\sqrt n\), normalization by \(\max\{R,K+\eta\}q\) gives the
all-plus certificate with
\(\gamma_{\rm tc}=K^2/(2(\beta+2KRB/3))\).  Its finite-confidence selected-library
form pays only the selected cap failures and centered matrix
Bernstein/Freedman tails.
A one-row sub-Weibull finite-confidence specialization chooses caps from
\(\log(A_*Mn/\delta)\), while the asymptotic cap-selection form chooses
\(b=sL^{1/\alpha}\) and
\(q=s\max\{L^{1/\alpha},\sqrt{\log(e\Delta)}\}\), pays only
\(J_mn_m e^{-L_m}\) for cap failures, and gives all-integer and sparse
active-menu schedules when \(L_m=c\log(e\Delta_m)\).
A finite-moment finite-confidence companion and its asymptotic operator-tail
schedule use \(U_m\)-Markov caps
\(b=sU_m^{1/p}\) and
\(q=s\max\{U_m^{1/p},\sqrt{\log(e\Delta)}U_m^{(2-p)_+/(2p)}\}\), giving the
same active schedules with polynomial cap-tail cost \(J_mn_mU_m^{-1}\).
A heterogeneous cap-profile form then lets each increment use its own cap
\(b_i\), with jump scale \(\max_i b_i\) but variance scale determined by the
aggregate truncated row/column second-moment profile; its finite-confidence
version pays only selected per-increment cap budgets and the selected matrix
tail.  Heterogeneous
sub-Weibull norm-tail scales use
\(q=\max\{\max_i s_iL^{1/\alpha},
(\log(e\Delta)n^{-1}\sum_i s_i^2)^{1/2}\}\).
A heterogeneous finite-moment selector similarly lets \(p_i,s_i,U_i\) vary,
using cap cost \(\sum_iU_i^{-1}\) and variance energy
\(\sum_i s_i^2U_i^{(2-p_i)_+/p_i}\) in the same profile; its
finite-confidence form pays the selected Markov cap budgets and the selected
matrix tail.
A fully tail-profile version, also with a finite-confidence form, accepts
deterministic survival profiles \(\pi_i\), with cap cost
\(\sum_i\pi_i(u_i)\) and variance energy
\(2\sum_i s_i^2\int_0^{u_i}t\pi_i(t)\,dt\), so exact distributional or
quantile bounds can be used directly without first replacing the increments by
a common tail class.
A quantile-budget selector then chooses
\(u_i=Q_i(\varepsilon\omega_i)\), charging only the stream budget
\(\varepsilon\) while allowing the weights \(\omega_i\) to be optimized across
heterogeneous increments.
A mixed sub-Weibull specialization makes this explicit for increment-dependent
tail constants and exponents through
\(b_i=s_i\log(eA_i/(\varepsilon\omega_i))^{1/\alpha_i}\) and variance energy
\(\sum_i(2A_i/\alpha_i)\Gamma(2/\alpha_i)s_i^2\), so the normalization keeps
average tail energy while the shares tune the worst cap.
When \(\alpha_i=\alpha\) is common, the optimal common cap is the smallest
\(r\) with \(\sum_iA_i\exp\{-(r/s_i)^\alpha\}\le\varepsilon/e\); the exact
capped variance energy is then
\(2\sum_i s_i^2\int_0^{r/s_i}t\min\{1,A_ie^{-t^\alpha}\}\,dt\).
The checkable envelope \(r=s_{\max}\log(e\sum_iA_i/\varepsilon)^{1/\alpha}\)
and \(V=(2/\alpha)\Gamma(2/\alpha)\sum_iA_is_i^2\) gives the same schedule,
and is the exact optimized cap when the scales are common.
A mixed finite-moment specialization similarly uses
\(\lambda_i=eA_i/(\varepsilon\omega_i)\), caps
\(b_i=s_i\lambda_i^{1/p_i}\), and moment energy
\(\sum_i\widehat\rho_i\), with
\(\widehat\rho_i=A_is_i^2\lambda_i^{(2-p_i)/p_i}\) for \(p_i<2\) and
\(\widehat\rho_i=A_i^{2/p_i}s_i^2\) for \(p_i\ge2\).
When \(p_i=p\) is common, the optimized shares are proportional to
\(A_is_i^p\); the resulting common cap is
\((e\sum_iA_is_i^p/\varepsilon)^{1/p}\), and the variance energy is
\((e/\varepsilon)^{(2-p)/p}(\sum_iA_is_i^p)^{2/p}\) for \(p<2\), or
\(\sum_iA_i^{2/p}s_i^2\) for \(p\ge2\).
Envelope schedules replace these exact aggregates by deterministic bounds
\(\bar S\ge\sum_iA_is_i^p\) and, for \(p\ge2\),
\(\bar W\ge\sum_iA_i^{2/p}s_i^2\), giving explicit power-growth
normalizations when \(\bar S,\bar W\) are polynomial in \(e\Delta\).
A variance-profile route
via noncommutative Khintchine gives \(O(\sqrt n)\) on every realization with
\(\|\sum_i A_i^2\|_{\mathrm{op}}=O(n/\log(eD))\).  Matrix Bernstein verifies this
profile for independent random contractions with
\(\|\sum_i\mathbb E A_i^2\|_{\mathrm{op}}=O(n/\log(eD))\), including
Haar-random rotations of fixed or random spectra with average Frobenius-square
\(O(D/\log(eD))\), and more generally conditional square-twirl and exact
finite-design rotations, including finite irreducible unitary-group rotations
and signed-permutation/monomial twirls, under the same average
Frobenius-square profile.  A block-design twirl theorem covers partially
randomized bases and random orthogonal decompositions by replacing the scalar
energy with the block-energy operator
\(\|\sum_{i,\alpha}d_{i,\alpha}^{-1}
\operatorname{tr}(P_{i,\alpha}T_i^2P_{i,\alpha})P_{i,\alpha}\|_{\mathrm{op}}\).
Approximate scalar and block designs are covered by the dominated-twirl
variant, which adds only the aggregate positive semidefinite twirl-error
profile to the same square-function budget.  A common-sign self-adjoint
square-profile menu version direct-sums active streams, pays
\(\log(eJ_mD_m)\), and removes the \(J_m\) factor from a sampled shared sign
search.  Its rank-sensitive version replaces \(J_mD_m\) by the certified total
active square-function rank \(R_m\), so low-rank active libraries pay
\(\log(eR_m)\).  Hermitian upper-triangular
independent-entry streams are controlled by their maximum cumulative row
variance, with whole-matrix cap tails giving Wigner/GOE/GUE-type unbounded
variants and one sign per whole matrix.  A Hermitian-dilation addendum gives
the same square-profile route for rectangular or non-self-adjoint random
operators from the left and right profiles
\(\|\sum_iX_iX_i^*\|_{\mathrm{op}}\) and
\(\|\sum_iX_i^*X_i\|_{\mathrm{op}}\), with predictable and active-menu
Borel--Cantelli versions; for centered conditionally independent-entry
rectangular matrices this is checked by maximum cumulative row and column
variance, with one sign per whole matrix, and whole-matrix cap tails give
sub-Weibull or finite-moment unbounded variants.  Its low-overlap specialization
gives checkable rectangular sparse-support criteria: row and column weighted
occupancies control the two profiles.  A realized-profile common-sign menu
version direct-sums the active rectangular library, pays
\(\log(eJ_m(d_{m,1}+d_{m,2}))\), and removes the \(J_m\) factor from a sampled
shared sign search.  Its rank-sensitive form replaces the ambient active
rectangular dimension by the certified total active left-plus-right
square-profile rank.  The dependent overlap-budget active-menu version only
needs summable certification of the realized left/right overlap budgets; a
tail-capped dependent version adds sub-Weibull or finite-moment coefficients by
paying cap-tail summability plus the capped realized overlap certificate.  A
support-rank common-sign overlap version replaces \(J_m(d_{m,1}+d_{m,2})\) by
the total active row-plus-column support-envelope rank, and the same rank
accounting survives tail capping under summable cap failures.  A coordinate-window
Bernstein version checks the row/column occupancy budgets at this support rank;
its predictable Freedman version allows adaptive row/column support exposure and
adds explicit arbitrary-growth, all-support-rank, and sparse-support-rank
schedules.
The parallel subspace-envelope Bernstein version checks random row/column
projection budgets inside deterministic low-rank envelopes.  Its predictable
Freedman analogue allows adaptive row/column projection supports under
conditional overlap budgets and has matching explicit support-rank schedules.
Scalar Bernstein verifies inhomogeneous bipartite support models.  A
rectangular random-subspace Bernstein variant replaces row/column
occupancies by the expected row and column projection sums, with Haar ranks
entering through \(d_1^{-1}\sum_i\rho_i^2r_i^L\) and
\(d_2^{-1}\sum_i\rho_i^2r_i^R\); both support variants have explicit
active-menu schedules paying \(J_m(d_{m,1}+d_{m,2})\), with
\(c_0a>\Gamma+2\) on all integer rectangular dimensions,
\(c_0a>\Gamma+1\) on sparse schedules, and sampled-certificate costs
\(J_m2^{-M_m}\).  A tail-capped coordinate-support form pays deterministic
operator-norm cap tails plus the capped row/column occupancy series, with a
cap-selection schedule for sub-Weibull and finite-moment sparse rectangular
blocks.  A tail-capped random-subspace form pays operator-norm cap exceptions
for unbounded low-rank or block coefficients; its cap-selection schedule has
the all-integer \(R>\Gamma+1\) and sparse \(R>\Gamma\) cap-tail thresholds.
Its rank-one dyad
specialization uses only the two expected
weighted frame operators, covering Haar-isotropic dyads, randomized feature
outer products, and nonuniform sampled rectangular atom dictionaries, again with
active-menu schedules for countable libraries.  Tail-capped dyad variants add
Gaussian, sub-Gaussian, sub-exponential, and heavy-tailed scalar weights by
paying the cap-exception probabilities plus the capped frame-Bernstein term;
their cap-selection schedule gives explicit sub-Weibull and finite-moment
scalar-weight criteria with the same all-integer and sparse cap-tail thresholds,
and its homogeneous isotropic specialization gives closed-form
\(R>\Gamma+1,\ c_0a>\Gamma+2\) all-integer and
\(R>\Gamma,\ c_0a>\Gamma+1\) sparse schedules.  A coordinate-sampled dyad
specialization replaces the frame budget by the largest row or column marginal
load, giving the same schedule for iid sampled-entry, random bipartite-edge,
and sparse observation matrices; its predictable form covers adaptive sampling,
without-replacement exposure, and permutation or matching windows from
one-step marginal-load envelopes, while a profile-capped variant keeps the full
row and column load arrays for nonuniform coordinate laws and varying caps,
with finite-menu all-integer/sparse schedules and sampled certificates; an
entry-intensity specialization checks adaptive inhomogeneous sparse matrices
from predictable weighted row/column degrees, with a closed-form cap schedule
for sub-Weibull or finite-moment entry weights and a position-dependent
cap profile charging rare heavy coordinates by \(q_{i,r,s}^2p_{i,r,s}\).
A matrix-Freedman version covers adapted dependent
random contractions whose predictable square profile is \(O(n/\log(eD))\); its
rank-sensitive common-corner form replaces \(D\) by the joint support rank
\(\kappa\) in both the logarithm and the Freedman dimension factor, and the
corresponding active-menu schedule pays \(\sum_h\kappa_{m,h}\).  A
two-scale square-function certificate now handles a rank-\(\kappa\) core plus a
full ambient residual under one signing; its predictable menu schedule pays
\(\sum_h\kappa_{m,h}\) for the low-rank core and \(J_mD_m\) only for the
residual.  A finite-depth multiscale version handles several low-rank layers
plus one ambient residual, with active structural cost
\(\sum_{h,q}\kappa_{m,h,q}+J_mD_m\).  A summable-depth version permits
countably many low-rank layers under a summable \(\sum_q\sqrt{a_q}\) budget,
with active menus paying the corresponding \(\sum_{m,h,q}\) structural tail.
A flag-compressible specialization generates those countable layers from
nested initial projection flags, covering multiresolution, wavelet,
spectral-window, and hierarchical basis streams under summable shell square
profiles.  A compression-error flag version replaces shell-square estimates by
predictable approximation budgets
\(\sum_i(\tau_{i,q-1}+\tau_{i,q})^2\), making spectral truncation and low-rank
tail estimates directly checkable.  A tail-envelope specialization gives
explicit common-flag rate criteria: polynomial tails beat
\(\exp(\gamma q^s)\) flag-rank growth when \(p>1+s/2\), and any positive power
tail works for geometrically growing rank windows.
A
rank-sensitive dominated square-twirl corollary gives the same replacement for
Haar, finite-design, signed-permutation, and approximate random-basis rotations
inside a common random low-rank corner; the matching finite-menu version makes
active low-rank random-basis libraries simultaneous at cost
\(\sum_h\kappa_{m,h}\).  A rank-sensitive predictable dominated-twirl version
gives the same low-rank replacement for online adaptive random-basis streams
chosen from the past.  Its
finite-menu schedule gives simultaneous almost-sure signings for active
libraries of arbitrary adapted square-profile streams after replacing \(D_m\)
by \(J_mD_m\), with common sampled batches paying only \(J_m2^{-M_m}\).  Its
predictable dominated-twirl specialization permits online adaptive scalar,
block, and approximate random-basis samplers with only the one-step PSD twirl
error added to the same budget.  The corresponding Borel--Cantelli array theorem
gives almost-sure eventual signings and sampled certificates whenever the row
budget tails and \(2^{-M_m}\) search losses are summable; explicit schedule
forms give all integer dimensions once \(n(D)\gtrsim\log^2(eD)\), and sparse
stretched-exponential dimension schedules under the one-row logarithmic-square
threshold.  A finite-menu predictable twirl schedule gives simultaneous
signings for active families of online random-basis samplers after replacing
the one-row structural count \(D_m\) by \(J_mD_m\), with common sampled batches
paying only \(J_m2^{-M_m}\).  Independent random sign search finds a certified
signing under the
same profile event with only an additional \(2^{-m}\) failure after \(m\) trials.
A
subspace-localized specialization covers adaptive coordinate supports, online
random subspaces, and conditionally sampled localized Wigner/GOE/GUE blocks
under the same predictable overlap scale.  Its finite-menu schedule gives
simultaneous signings for active adaptive support/subspace libraries after
replacing \(D_m\) by \(J_mD_m\), with common sampled batches paying only
\(J_m2^{-M_m}\).  A
conditional scalar-multiplier corollary covers arbitrary
deterministic or random self-adjoint contraction directions after independent
bounded Bernoulli thinning or scalar amplitudes with conditional second moments
\(O(1/\log(eD))\), and a predictable version covers adaptive online directions
without conditional independence; its triangular-array schedule form gives the
same all-integer \(n(D)\gtrsim\log^2(eD)\) and sparse stretched-exponential
almost-sure regimes under the half-budget conditional second-moment bound.
Finite menus of such adaptive scalar-multiplier streams are simultaneous after
replacing \(D_m\) by \(|\mathcal H_m|D_m\), with common sampled batches paying
only \(|\mathcal H_m|2^{-M_m}\).  Exhaustive countable scalar-multiplier menus
are controlled on prefixes \(h\le J_m\) at cost \(J_mD_m\), so every fixed
stream is eventually covered.  A
hybrid centered-noise version preserves any already-known
\(K\sqrt n\) core signing, including algebraic or finite-lift core signings,
after adding independent bounded scalar noise in arbitrary matrix directions,
with an explicit matrix-Bernstein tail; a martingale-difference version removes
that conditional independence and allows bounded matrix-valued martingale noise,
and the predictable matrix-noise branch has a Borel--Cantelli schedule theorem
for all integer dimensions and stretched-exponential sparse dimensions whenever
the predictable quadratic-variation budget is \(O(n/\log(eD))\);
its active-prefix countable-menu version pays \(J_mD_m\), so every fixed
already-signed core-plus-noise model in a countable library is eventually
covered;
its sampled localized specialization gives the same almost-sure schedules for
adaptive coordinate-support and conditional Haar-subspace matrix-noise arrays
under the predictable weighted-overlap or average-rank budget;
the localized branch also has an active-prefix countable-library form at the
same \(J_mD_m\) cost, using weighted-overlap or Haar average-rank budgets
directly;
a scalar variance-domination corollary gives the same signed-core tail whenever
\(\mathbb E(Y_i^2\mid\mathcal F_{i-1})\preceq\theta_i\rho_i^2I\) with
independent random variance intensities, and an anisotropic matrix-profile
version allows random PSD variance profiles with a matrix-Bernstein budget tail;
a trace-budget variant removes the ambient-dimension factor from the budget
exception when the total random variance trace is controlled, and a
trace-envelope/low-rank specialization makes this explicit for random
low-rank, spiked, and effective-dimension variance profiles; the corresponding
Schatten-envelope corollary interpolates between trace, Frobenius-plus-rank,
and operator-plus-rank control, and a sub-gamma trace-envelope variant allows
unbounded envelopes with a Bernstein-type total-envelope tail, while a
quantile/finite-moment trace-envelope variant covers heavy-tailed total trace
budgets without independence, and its summable Borel--Cantelli form gives
eventual Spencer-scale bounds for triangular arrays with trace-budget quantiles
at \(n/\log(eD)\); normalized \(r\)-th moment summability of those heavy-tailed
budgets is enough to check this criterion, and a termwise moment version checks
random spike, low-rank, trace-class, and Schatten atoms without independence;
a termwise-tail version only needs summable individual truncation failures, and
a scale-mixture specialization handles deterministic atom weights with
heavy-tailed multipliers via aggregate tail summability; when all atoms share
one multiplier, a common-scale version removes the atom-count union factor,
and finite or countable grouped-scale versions pay one tail per volatility
group.  An optimized tail-profile form allows exact survival functions and
log-Weibull/lognormal tails, while the heterogeneous polynomial and sub-Weibull
specializations have optimized share-cost criteria.  A common-profile aggregate
criterion handles shared tail shapes by proportional mass shares, and an
optimized profile-menu version aggregates mixed tail families by best
class-level trace shares.  A finite heterogeneous polynomial-index menu is
controlled by the worst tail exponent \(\alpha_0\), with the same
\(\alpha_0 b>p+\alpha_0 u+(\alpha_0+1)w+1\) count/trace rule, while a
countable weighted heterogeneous power-decay version gives
\(\alpha_0 b>p+\alpha_0 u+1\) under summable class-share/trace decay.  A finite
heterogeneous stretched-exponential profile menu has logarithmic budget rule
\(\delta>p+w+1\), with \(\delta=c_0a^{\beta_0}\), after paying the active
class-count exponent; its countable weighted trace-decay version removes this
penalty and gives \(\delta>p+1\).  The analogous finite log-Weibull menu has
the same exponent rule with a subpolynomial budget margin, and the countable
weighted log-Weibull version removes the count penalty once the residual
class-index log-Weibull series is summable.  In the
common-index polynomial case there is a closed-form effective class mass and
schedule rule
\(\alpha b>(\alpha+1)p+1\), with logarithmic endpoint
\((\alpha+1)q-\alpha c<-1\).  A finite class-count version reduces to
\(\alpha b>p+\alpha u+(\alpha+1)w+1\), and a countable power-decaying class
version reduces to \(\alpha b>p+\alpha u+1\) when
\(r+\alpha s>\alpha+1\).
A stretched-exponential schedule has \(ca^\beta>p+1\) and endpoint
\(q-ca^\beta b<-1\), plus a log-Weibull schedule for subpolynomial budget
ratios; a regularly varying log-corrected schedule has exponent rule
\(\alpha a>p+1\) and endpoint \(q+\gamma-\alpha b<-1\); the group-level
common-index polynomial case has its own explicit effective tail mass.
The sub-gamma hypothesis is
verified termwise from independent sub-gamma envelopes and specialized to
atomized and countable bounded, gamma/chi-square, gamma-atom, countable gamma-atom,
weighted-chi-square/real-Wishart,
weighted-exponential/complex-Wishart, noncentral/correlated Gaussian/Wishart,
compound-Poisson, and mixed-atom trace
envelopes when available;
the scalar specialization allows predictable directions with the matching
matrix-Freedman tail.  Localized specializations handle online
coordinate-supported or subspace-localized matrix/scalar martingale noise through
realized predictable overlap, conditional predictable overlap, and
weighted-overlap operators.
The signed-core matrix-noise theorem is also specialized to adaptively sampled
coordinate supports with coordinate-dependent predictable inclusion bounds and
conditional Haar subspaces via random weighted conditional inclusion-probability
and average-rank budgets, with an explicit budget-tail term.
Predictable amplitude weights sharpen the conditional-overlap parameter from
\(\sum_i\mathbb E(P_i\mid\mathcal{F}_{i-1})\) to
\(\sum_i\rho_i^2\mathbb E(P_i\mid\mathcal{F}_{i-1})\).
For coordinate supports with deterministic or independent random inclusion
profiles, and for Haar subspaces with independent random amplitudes, scalar
Bernstein makes the weighted-budget exception explicit; in the homogeneous
coordinate case the overlap scale becomes \(p\sum_i\mathbb E\rho_i^2\), while
fixed orthogonal block decompositions have the same tail with the coordinate
union factor replaced by the number of blocks, and
fixed or random Haar ranks use the scalar budget
\(\sum_i\mathbb E[(R_i/D)\rho_i^2]\).  The Haar assumption may be replaced by
conditional isotropy, or by the weaker spectral domination
\(\mathbb E(P_i\mid\mathcal F_{i-1})\preceq\alpha_iI\), with random intensity
\(\alpha_i\).  Mixed
coordinate-or-Haar samples are
covered by the same profile after replacing the coordinate inclusion bound by
\(M_iQ_{ij}+(1-M_i)R_i/D\).  A random-amplitude localized-noise menu converts
all these scalar Bernstein budget tails into Borel--Cantelli all-integer and
sparse dimension schedules whenever the budget exceptions and matrix-Freedman
row tails are summable; in the homogeneous subcases this reduces to the
effective-energy rule \(\beta\sum_i\mathbb E\rho_i^2=O(n/\log(eD))\), with
\(\sum_i\mathbb E\rho_i^2\gtrsim\log(eD)\) only needed to make the scalar
budget exception summable after the relevant coordinate/block union count.
The countable-menu version controls active prefixes by replacing \(N_m\) with
\(\sum_{h\le J_m}N_{m,h}\), so every fixed localized random-amplitude model is
eventually covered.
A common-sign version says that when one core signing balances every
active-prefix core, that same signing survives all active-prefix localized
random-amplitude perturbations under the same union-count payment.
A finite-lift common-menu corollary supplies such shared core signings
deterministically whenever the active-prefix direct-sum lift entropy is \(O(n)\).
Consequently finite-lift core libraries absorb localized random-amplitude noise
under the same active-prefix union-count and homogeneous energy schedules.
Its homogeneous active-prefix specialization pays
\(\Lambda_m^{(J)}=\sum_{h\le J_m}N_{m,h}\), giving all-integer energy margin
\(R>8(\Gamma+1)/3\) under \(\Lambda_D^{(J)}\le(eD)^\Gamma\) and sparse margin
\(R>8\Gamma/3\).
The homogeneous energy rule also has the same common-sign form whenever the
active-prefix cores share one signing.
It also covers random localized,
projection, frame, and sparse-support ensembles whenever their weighted empirical
overlap or weighted frame-operator norm is \(O(n/\log(eD))\).  An active-menu
dependent version only needs summable certification of realized overlap budgets,
and a tail-capped dependent version adds sub-Weibull or finite-moment localized
blocks by paying cap-tail summability plus the capped overlap certificate.  A
common-sign version direct-sums the active menu and gives one signing for all
dependent localized streams after replacing \(\log(eD)\) by
\(\log(eJ_mD_m)\).  A support-rank common-sign variant replaces \(J_mD_m\)
by the total active support-envelope rank for localized libraries, including
tail-capped unbounded versions with summable cap failures.  A coordinate-window
Bernstein version checks the active support-rank budgets from scalar occupancy
tails, while a predictable Freedman version allows adaptive coordinate-support
exposure and gives explicit arbitrary-growth, all-support-rank, and
sparse-support-rank schedules.  A subspace-envelope Bernstein version does the
same for random projection supports inside deterministic low-rank envelopes.
The predictable subspace-envelope Freedman version removes independence for
adaptive projection supports with controlled conditional overlap and has
matching explicit support-rank schedules.  For independent inhomogeneous
coordinate supports, a scalar Bernstein union bound
verifies this event from the weighted occupancies \(\max_j\sum_i\rho_i^2{\bf 1}_{\{j\in S_i\}}\);
the old support-density \(O(1/\log(eD))\) theorem is the homogeneous
specialization.  For independent random subspace projections, matrix Bernstein
gives the same conclusion when
\(\|\sum_i\rho_i^2\mathbb E P_i\|_{\mathrm{op}}=O(n/\log(eD))\), including Haar
weighted-rank condition \(D^{-1}\sum_i\rho_i^2r_i=O(n/\log(eD))\).  Independent
random amplitudes are also covered after replacing these budgets by
\(\max_j\sum_i p_{ij}\mathbb E\rho_i^2\) and
\(\|\sum_i\mathbb E(\rho_i^2P_i)\|_{\mathrm{op}}\), with an explicit
random-amplitude schedule corollary giving the same all-integer, sparse,
arbitrary-growth, and sampled-certificate thresholds.  Fixed orthogonal
block-support ensembles pay the block count \(B\), not the ambient coordinate
count \(D\), with budget \(\max_\alpha\sum_iq_{i\alpha}\mathbb E\rho_i^2\) and
schedule margin \(\log(eB_m)\).  Random conditional inclusion profiles are
also allowed, replacing deterministic budgets by
\(\max_j\sum_i\mathbb E(\rho_i^2Q_{ij})\) or
\(\max_\alpha\sum_i\mathbb E(\rho_i^2Q_{i\alpha})\), which covers random support
sizes even when they are correlated with amplitudes.  The random-profile route
now has its own variance-sensitive Borel--Cantelli array form: with
\(N_m=D_m\) or \(B_m\), summable
\(N_m\exp\{-t_m^2/(2(\sigma_m^2+t_m/3))\}\) gives eventual profile signings and
sampled certificates.  A predictable profile/Freedman variant also permits
online coordinate or block sampling from \(\mathcal F_{i-1}\)-measurable
profiles, paying only the realized predictable-overlap budget exception; its
schedule form has the same \(\log(eN_m)\) arbitrary-growth rule and the same
all-integer/sparse thresholds, conditional on summable budget exceptions.
Finite menus of such predictable coordinate/block support models are
simultaneous after replacing \(N_m\) by \(L_m=\sum_hN_{m,h}\) and summing the
budget exceptions.  Exhaustive countable menus are diagonalized by controlling
prefixes \(h\le J_m\) at cost \(L_m^{(J)}=\sum_{h\le J_m}N_{m,h}\), so every
fixed online support model is eventually covered.  A
deterministic-envelope specialization has no budget exception: fully adaptive
supports with conditional inclusion bounds \(p_{ij}\) or \(q_{i\alpha}\) and
deterministic amplitude caps use budgets
\(\max_j\sum_i a_i^2p_{ij}\) or
\(\max_\alpha\sum_i a_i^2q_{i\alpha}\).  The deterministic-envelope schedule
form keeps the same \(\log(eN_m)\) arbitrary-growth, all-integer, sparse, and
sampled-certificate thresholds with no budget-exception summability term.
Finite menus of such coordinate/block envelope models are simultaneous after
replacing \(N_m\) by \(L_m=\sum_hN_{m,h}\), with common sampled batches paying
only \(|\mathcal H_m|2^{-M_m}\).  Exhaustive countable deterministic-envelope
menus control prefixes \(h\le J_m\) at cost
\(L_m^{(J)}=\sum_{h\le J_m}N_{m,h}\), with no budget-exception series.  A
Borel--Cantelli triangular-array corollary makes these
coordinate and subspace conclusions eventual almost-sure statements under the
corresponding summable Bernstein failure series; under the half-budget condition
this is automatic
along all integer dimensions from \(n_D\ge a\log^2(eD)\), \(a>32/(3C_0)\), and
along sparse stretched-exponential dimension schedules from \(a>16/(3C_0)\).
The same weighted support families have
sampled-sign certificates with only an additional \(2^{-M}\) finite-search
failure; logarithmically many trials give eventual sampled certificates:
\((1+\delta)\log_2(eD)\) trials along all integer dimensions and
\((1+\delta)\log_2(e+m)\) trials along sparse schedules.  More generally,
\(n_m\ge(16/(3C_0))\log(eD_m)(\log(eD_m)+(1+\gamma)\log(e+m))\) makes the
weighted-support and sampled-search failures summable for arbitrary triangular
arrays.  The
commutative AF specialization is now explicit: a clopen
partition with \(M\) atoms and oscillation \(\eta\) gives discrepancy
\(C\sqrt{n\log(e+M/n)}+2n\eta\), while interval-range functions force
\(\Omega(\eta^{-1})\) atoms for any finite-stage certificate at scale \(\eta\),
hence \(\Omega(\sqrt n)\) atoms at the Spencer approximation scale.  A metric
version bounds the discrepancy by a common modulus of continuity and the clopen
partition growth function \(M_X(r)\); for \((L,\alpha)\)-Hölder tuples with
\(M_X(r)\le Ar^{-d}\) it gives \(O(\sqrt{n\log n})\), and the Spencer scale
\(O(\sqrt n)\) follows when \(d\le2\alpha\).  Equivalently, for compact
zero-dimensional metric spaces, finite upper clopen Minkowski dimension gives
the \(O(\sqrt{n\log n})\) metric AF bound, and
\(\overline{\dim}_{\rm cl}X<2\alpha\) gives Spencer scale for Hölder tuples;
the critical endpoint \(\overline{\dim}_{\rm cl}X=2\alpha\) is also
Spencer-scale whenever the upper clopen \(2\alpha\)-Minkowski content is finite.
At the logarithmic critical window
\(M_X(r)\lesssim r^{-2\alpha}(\log(e/r))^\beta\), the bound is still
\(O(\sqrt n)\) for \(\beta\le0\) and improves the generic logarithmic loss to
\(O(\sqrt{n\log\log n})\) for \(\beta>0\).
The clopen-growth input can now be replaced by ordinary metric covering
numbers: every finite ball cover of a compact zero-dimensional space has a
clopen shrinking with the same cardinality, so \(N_X(\rho)\lesssim\rho^{-d}\)
implies the same H\"older threshold \(d\le2\alpha\), and the same logarithmic
critical-window estimates hold with \(N_X\) in place of \(M_X\).
Equivalently, finite ordinary upper metric Minkowski dimension gives the
\(O(\sqrt{n\log n})\) bound, ordinary dimension \(<2\alpha\) gives Spencer
scale, and finite ordinary upper metric \(2\alpha\)-Minkowski content gives the
critical endpoint.
The same package is stated in Assouad/doubling language: upper Assouad growth
\((A,s)\) gives the identical \(s\le2\alpha\) threshold, and a \(D\)-doubling
zero-dimensional model is Spencer-scale whenever \(\log_2D\le2\alpha\).
Finite unions keep the largest ordinary covering exponent and the largest
ordinary upper metric dimension, with the same largest-exponent rule in
Assouad/doubling language; finite max-metric products add ordinary covering
exponents, ordinary Minkowski dimensions, critical contents, and
Assouad/doubling exponents, so the product threshold is governed by the sum of
the factor exponents.  Combining the two rules, finite product-stratified
unions are governed by the maximum, over strata, of the sum of the factor
exponents in that stratum.
Zero-dimensional H\"older resolutions transfer the same entropy bound to
non-AF targets.  Every compact metric space has scale-flexible zero-dimensional
Lipschitz resolutions: for \(0<\theta<1\), the profile is preserved as
\(N_X(s)\le N_Y(\theta s)\) with Lipschitz loss \(2/(1-\theta)\).  This gives
the direct bound
\(C\sqrt{n\log(e+N_Y(\rho)/n)}+2n\omega(16\rho)\) by \(\theta=1/2\).
Consequently every fixed compact target has qualitative \(o(n)\) discrepancy
for any sequence of contraction \(n\)-tuples sharing one equicontinuity modulus.
Equivalently, every relatively compact class of real contraction test functions
has hereditary discrepancy \(o(n)\).
In fact, the qualitative principle is Banach-space level:
every totally bounded subset of the unit ball of a normed space satisfies
\(\Delta_n\le N_{\mathcal K}(\eta)+n\eta\), hence hereditary \(o(n)\);
polynomial Banach entropy \(N_{\mathcal K}(\eta)\lesssim\eta^{-d}\) gives
\(\Delta_n=O(n^{d/(d+1)})\), pure stretched-exponential Banach entropy gives
\(\Delta_n=O(n/(\log n)^{1/q})\), and signed-log stretched entropy
\(\exp(A\eta^{-q}(\log(e/\eta))^\gamma)\) gives the refined rate
\(n(\log n)^{-1/q}(\log\log n)^{\gamma/q}\).  Thus compact operator images of
unit balls have sublinear signed sums, with rates governed by their image
entropy.
Direct dyadic entropy-number data now gives the sharper scale-by-scale package
\(\Delta_n\le2^{k-1}+n e_k^{\rm int}\): polynomial decay
\(e_k^{\rm int}\lesssim k^{-s}\) yields
\(\Delta_n\lesssim n(\log n)^{-s}\), log-refined decay adds the expected
\((\log\log n)^\beta\) factor, and \(e_k^{\rm int}\lesssim e^{-bk}\) gives
\(\Delta_n\lesssim n^{\log2/(b+\log2)}\).  The package now converts standard
external entropy numbers to internal ones at only a factor \(2\), and records
the finite-union and finite-Minkowski-sum rules for carrying these rates through
compact class decompositions.  It also handles exponentially weighted summable
Minkowski series: uniform \(e_k\lesssim k^{-s}(\log k)^\beta\) component
entropy gives \(e_m\lesssim m^{-s}(\log m)^{s+\beta_+}\) and hence
\(\Delta_n\lesssim n(\log n)^{-s}(\log\log n)^{s+\beta_+}\), while exponential
component entropy gives \(\Delta_n\lesssim n\exp(-c\sqrt{\log n})\).
The weighted-series statement now has a profile form:
\[
  e_m\lesssim \inf_{r\le m}\{S_r\varphi(1+\lfloor(m-1)/r\rfloor)+T_r\},
\]
with the analogous dyadic \(\Delta_n\) bound, where
\(S_r=\sum_{j\le r}a_jb_j\) when
\(e_k(\mathcal K_j)\lesssim b_j\varphi(k)\), and \(T_r\) is the remaining
radius tail.
The same profile is optimized abstractly: if
\(S_r\lesssim r^\delta(\log r)^\eta\),
\(T_r\lesssim r^{-\gamma}(\log r)^{-\tau}\), and
\(\varphi(k)\lesssim k^{-s}(\log k)^\beta\), then
\[
  e_m\lesssim
  m^{-s\gamma/(s+\delta+\gamma)}
  (\log m)^{(\gamma(\beta+\eta)-\tau(s+\delta))/(s+\delta+\gamma)}.
\]
Thus polynomial weights with component amplitudes
\(j^\rho(\log j)^\eta\), \(\rho>u-1\), yield
\[
  e_m\lesssim
  m^{-s(u-1)/(s+\rho)}
  (\log m)^{((u-1)(\beta+\eta)-v(s+\rho))/(s+\rho)}.
\]
More generally, stretched-exponential weights \(a_j\lesssim e^{-cj^\theta}\)
replace the finite polynomial entropy log-power by \(s/\theta+\beta_+\), and
turn exponential component entropy into
\(\Delta_n\lesssim n\exp(-c(\log n)^{\theta/(1+\theta)})\).
For polynomial weights \(a_j\lesssim j^{-u}\), \(u>1\), the optimized exponent
is \(\sigma=s(u-1)/(s+u-1)\); exponentially entropic summands are then limited
by the tail scale \(u-1\), up to logarithmic powers.
A log-balanced refinement covers
\(a_j\lesssim j^{-u}(\log(e+j))^{-v}\): if component entropy is
\(e_k\lesssim k^{-s}(\log k)^\beta\), then
\[
  e_m\lesssim m^{-\sigma}(\log m)^\lambda,
  \qquad
  \lambda=\frac{(u-1)\beta-sv}{s+u-1},
\]
with the matching dyadic
\(n(\log n)^{-\sigma}(\log\log n)^\lambda\) discrepancy rate; exponentially
entropic summands give the tail-limited log power \(u-1-v\).
H\"older images are now included as well:
an \(\alpha\)-H\"older parametrization multiplies the entropy exponent and
log-power by \(\alpha\), so parametrized compact classes inherit the matching
logarithmic discrepancy rates.
The Banach section now also records a finite-rank approximation envelope:
if a compact image is within \(a_m\) of an \(m\)-dimensional subspace, then
\[
  \Delta_n\le 3(1+6/\eta)^m+n(\eta+a_m),
\]
so approximation-number decay \(a_m\lesssim m^{-s}\) gives
\(\Delta_n\lesssim n(\log\log n/\log n)^s\), log-refined decay
\(a_m\lesssim m^{-s}(\log m)^\beta\) gives the same rate multiplied by
\((\log\log n)^\beta\), while
\(a_m\lesssim\exp(-bm^\theta)\) gives
\(\Delta_n\lesssim n\exp(-c(\log n)^{\theta/(1+\theta)})\), including
\(\exp(-c\sqrt{\log n})\) for \(\theta=1\).
Hilbert singular-value decay gives the same rates; Schatten \(p\)-class gives
exponent \(1/p\).

For compact metric targets with \(N_Y(\rho)\le A\rho^{-d}\), the loss is
\[
  O\!\left(\sqrt{n\log(e+n^{d/(2\alpha)-1})}\right),
\]
and it is Spencer-scale when \(d\le2\alpha\); the logarithmic critical window
and Minkowski/content thresholds also hold without zero-dimensionality.  Local
Assouad/doubling
profiles also pass to arbitrary compact targets: upper Assouad exponent
\(s\le2\alpha\), in particular \(\log_2D\le2\alpha\) for compact
\(D\)-doubling targets, gives Spencer scale, and signed logarithmic Assouad
profiles have the same critical-window split.  Beyond polynomial entropy,
stretched-exponential compact targets with
\(\log N_Y(\rho)\lesssim\rho^{-q}(\log(e/\rho))^\gamma\) still have sublinear
H\"older discrepancy
\[
  O\!\left(n^{(\alpha+q)/(2\alpha+q)}
    (\log n)^{\alpha\gamma/(2\alpha+q)}\right).
\]
H\"older quotient maps divide
the model exponent by the quotient exponent: if \(Y\) is a
\(\beta\)-H\"older image of a compact model with covering exponent \(s\), the
threshold is \(s\le2\alpha\beta\), and a \(\beta\)-H\"older quotient of a
compact \(D\)-doubling model is Spencer-scale when
\(\log_2D\le2\alpha\beta\).  The quotient version also keeps the signed
logarithmic critical window: at \(s=2\alpha\beta\), nonpositive log power remains
Spencer-scale and positive log power costs only \(\sqrt{n\log\log n}\).  For
snowflaked metrics \(d^\vartheta\), this becomes \(s\le2\alpha\vartheta\),
with the same signed logarithmic endpoint at \(s=2\alpha\vartheta\).
Finite unions of arbitrary
compact metric targets keep the largest covering, Minkowski, and
Assouad/doubling exponent, while finite max-metric products add ordinary
covering exponents, Minkowski contents, and Assouad/doubling exponents after
the same resolution, so the Spencer threshold is governed by the maximum for
unions and the sum for products.  The same finite-union/product rules now track
signed logarithmic covering powers: only top polynomial exponents contribute to
the log correction, while products add the log powers.  Finite unions of
product-parametrized strata use the maximum over stratum sums, and H\"older
stratum parametrizations divide each stratum sum by the parametrization
exponent; at critical exponent \(2\alpha\), this gives Spencer scale for
nonpositive top log power and otherwise the
\(\sqrt{n\log\log n}\) window.  The binary Cantor
resolution gives \([0,1]^D\), and compact subsets for \(0<\alpha\le1\), the
Euclidean threshold \(D\le2\alpha\); compact targets with finite
\(\beta\)-H\"older \(D\)-atlases have the divided threshold
\(D/\beta\le2\alpha\), with compact Lipschitz manifolds and finite polyhedra as
the \(\beta=1\) case.  More generally, finite H\"older atlases whose chart
domains have profiles \(\rho^{-d_\ell}(\log(e/\rho))^{\gamma_\ell}\) inherit
the same critical-window rule with top divided exponent
\(\max_\ell d_\ell/\beta_\ell\) and top log power among the critical charts.
An exhaustive finite-scale atlas corollary now tracks one-sided Hausdorff stage
error \(\delta_r\) and stage entropy budgets \(A_r\): at the signing scale
\(\rho_n\asymp n^{-1/(2\alpha)}\), the loss is exactly
\(\sqrt{n\log(e+A_{r_n}n^{s/(2\alpha)-1}(\log n)^\gamma)}\), so growing stages
are allowed up to the unused entropy slack.  In particular, if
\(\delta_r\lesssim r^{-\eta}\) and \(A_r\lesssim
r^\lambda(\log r)^\kappa\), the effective exponent is
\(s+\lambda/\eta\) with log power \(\gamma+\kappa\).  If instead
\(\delta_r\lesssim e^{-cr^\theta}\) and \(A_r\lesssim
r^\lambda(\log r)^\kappa\), there is no polynomial exponent penalty: only the
log power changes to \(\gamma+\lambda/\theta\), with a secondary
\((\log\log n)^\kappa\) factor.
The ultrametric specialization gives the same bounds from ball-covering growth,
and for the \(q\)-symbol Cantor shift with scale \(\theta\) the threshold is
\(\alpha\ge\frac12\log q/\log(1/\theta)\); finite products add these
dimensions, giving the threshold
\(\alpha\ge\frac12\sum_j\log q_j/\log(1/\theta_j)\) for product symbolic
models with the max metric.

Recent additions record the tensor-reversal proof for arbitrary banded
multilevel block Hankel matrices, the sample-optimal norming-grid count and
formal anisotropic $\prod_j(2w_j+1)$ banded theorem, its matching anisotropic
circulant lower bound already in the self-adjoint/block-circulant subclass, the
self-adjoint product-bandwidth regime
$\prod_j(2w_j+1)\le n$ with unrestricted ambient dimension, and the fixed-level
full-bandwidth multilevel block Toeplitz/Hankel theorem via the
Carlsson--Perfekt symbol extension, with constants depending on \(r,b\).  Lean
companions:
\texttt{paper12/lean} machine-checks the entropy-obstruction inequality, the
one-dimensional and multilevel interior-index counts, the grid-norming
arithmetic, the tensor-reversal and reversal-symmetry index/matrix arithmetic,
the multilevel band constant, and the entropy-absorption/direct-sum bookkeeping
(not the analytic norm-equality step), with no \texttt{sorry}/\texttt{admit}.

\textbf{13. \texttt{Residual\_Finite\_Horizon\_Fusion\_v1.pdf}} (18pp, Thread A).
The finite-horizon layer of paper 1's Hilbert-cube problem.  A positive
residual theorem: for any countable family of dense-open window sequences there
is a comeager set in $\mathcal F(Q_\infty)=C(Q_\infty,Q_\infty)^{\mathbb N}$
whose every system admits finite-horizon windowed DC1 witnesses of every order,
at every tail, every rational proximal scale and tolerance, and the fixed
separation scale --- the engine is a finite-horizon stability lemma whose
tracking modulus does \emph{not} decay with coordinate depth, against the depth
obstruction of paper 1.  The residual set now also contains finite scrambled
packets of every prescribed cardinality: after binary coding across the
separation blocks, all points in the packet share the same high-density
proximal times while each pair has a positive-density separation set.  And a
sharp finite-certificate obstruction: every
open set certified by a bounded window of coordinate maps that lies in the
infinite-metric-mean-dimension (or scale-law entropy) class is empty, by an
explicit coordinate squeeze past the certifying window.  This rules out the
natural finite-horizon method class for the full all-tail conclusions of
Problem 7.1 without deciding whether those conclusions are comeager.  The same
squeeze also rules out a finite-horizon certificate for the all-tail symbolic
semifactor, but it does not rule out a bare uncountable DC1 scrambled set; the
obstruction applies to the growth/entropy-bundled conclusions, not to
cardinality-only DC1.

\textbf{14. \texttt{Random\_Toeplitz\_Discrepancy\_v1.pdf}} (215pp, Thread B).
Operator-norm discrepancy
\[
  \mathrm{disc}_{\mathrm{op}}(A_1,\dots,A_n)
  =\min_{\varepsilon\in\{\pm1\}^n}
  \left\|\sum_i\varepsilon_iA_i\right\|_{\mathrm{op}}
\]
for independent random banded Toeplitz/Hankel contractions.  When the
half-bandwidth $w\le n$ (and $d\ge4w$)
the deterministic banded theorem of paper 12 already gives $O(\sqrt n)$ on every
realization, with a constant free of the dimension $d$ and $w\mapsto w^r$ at
multilevel.  For full bandwidth the paper records three settled endpoints: the
deterministic bound \(C\sqrt{n\log(e+d/n)}\), which pathwise settles bounded
aspect ratio \(d/n=O(1)\); a fixed-sign theorem showing that the
expected-normalized full random Toeplitz/Hankel model already has
\(O(\sqrt n)\) discrepancy with high probability, every polynomial-sized
prescribed deterministic sign list, and every independent random sign list with
polynomial expected size, is simultaneously \(O(\sqrt n)\). More generally,
every finite sign family of size \(M\) is simultaneously controlled at scale
\[
  O_K\!\left(\sqrt{n\left(1+\frac{\log M}{\log d}\right)}\right),
\]
so the whole signing cube has scale
\(O_K(\sqrt{n(1+n/\log d)})\), and in the logarithmic sample window
\(n\le\lambda\log d\) every sign pattern is simultaneously
\(O_{\lambda}(\sqrt n)\).  A profile-normalized fixed-sign theorem extends this
deterministic sign control to independent non-identical lag variables with
arbitrary deterministic sub-Gaussian variance envelopes, normalized by
\((\sum_{|k|<d}\sigma_{i,k}^2)^{1/2}\sqrt{\log(e+d)}\), and its entropy form
controls every sign family of size \(M\) at scale
\(O_K(\sqrt{n(1+\log M/\log d)})\), with an almost-sure triangular-array
version under the summability condition \(\sum_m d_m^{-L_m}<\infty\) and
explicit polynomial/stretched-exponential dimension schedules for polynomial
sign families and whole-cube logarithmic windows.  A coefficient-family
entropy variant replaces signs by arbitrary finite coefficient sets with
\(\ell^2\)-radius \(R\), giving scale
\(O_K(R\sqrt{1+\log M/\log d})\), and the same Borel transfer gives eventual
\(O(R_m)\) control for coefficient catalogs of size at most \(d_m^\alpha\) in
the polynomial and stretched-exponential dimension schedules.  The
expected-normalized sampled robust-certificate package also transfers to actual
full Toeplitz/Hankel operator norm in the
unbounded-aspect range, because the same fixed-sign theorem with \(n=1\)
supplies the needed uniform \(O(1)\) bound on the normalized summands; and a
realization-normalized convex-concentration
theorem.  The key input is the empirical edge profile
\[
  \max_{\omega^d=1}\sum_i
  |P_i(\omega)|^2/\|T_d(\xi^{(i)})\|_{\mathrm{op}}^2
  =O(n/\log d)
\]
for samples with \(n\le d\), \(\log d=o(n)\), proved by splitting off rare
low-normalizer samples.  A normalizer-tail criterion shows that this profile,
and hence \(O(\sqrt n)\) discrepancy, follows for any sub-Gaussian law satisfying
\(\Pr\{\|T_d(\xi)\|_{\mathrm{op}}<a\,\mathbb E\|T_d(\xi)\|_{\mathrm{op}}\}
\,\min\{n,\log^2(e+d)\}=o(1)\); the same profile also gives scalar edge-grid sign moments
\((Ck\,n/\log d)^k\) for every \(2k\)-th moment and a sign-averaged edge
maximum \(O(\sqrt n)\).  The empirical core is now data-dependent: if \(M_a\)
is the realized low-normalizer count, then outside the standard \(d^{-A}\)
events the profile is \(O(n/\log d+M_a\log d)\), and the signing scale is
\(O(\sqrt{n+M_a\log^2 d})\).  In particular, when
\(M_a=O(n/\log^2 d)\), the profile, random-sign tail, sign-averaged maximum,
randomized-search amplification with success-count tails, deterministic certificate, logarithmic-radius
Hamming stability with sampled robust-center tails, and scalar edge-grid moment bounds all recover the
\(O(\sqrt n)\) scale.  A deterministic Borel transfer template upgrades every
eventual edge-grid robust-center count in the full-bandwidth range
\(d_m/n_m\to\infty\), \(\log d_m=o(n_m)\), to actual Toeplitz/Hankel
operator-norm robust certificates with the same centers and Hamming radii.  A
dyadic normalizer refinement now replaces the single
bad-count penalty by
\(W_a=\sum_r4^rM_{a,r}+\ell M_{a,\infty}\), \(\ell=\log(e+d)\), so samples
whose normalizers are only moderately low are charged by their dyadic depth;
the full profile/signing/moment package holds at scale
\(O(\sqrt{n+W_a\ell})\), and returns to \(O(\sqrt n)\) whenever
\(W_a=O(n/\ell)\).  The same dyadic budget now has a one-matrix tail
criterion: if \(\mu_{d,a}^{\rm dy}=\mathbb E W_a/n\), Markov gives matrix
failure \(O(d^{-A})+O(\ell\mu_{d,a}^{\rm dy})\), while the large-sample
condition \(\mu_{d,a}^{\rm dy}\le\theta/\ell\) gives the sharper Bernstein
cost \(\exp\{-c n/\ell^2\}\).  An inverse-moment criterion further bounds
\(\mu_{d,a}^{\rm dy}\) by the truncated mass
\(\mathbb E[1_{\{\|T_d\|<aN_d\}}\min\{(aN_d/\|T_d\|)^2,\ell\}]\), and by
inverse \(p\)-moments with an explicit logarithmic loss; its Borel form gives
eventual profiles, sampled robust certificates, robust-fraction tests,
polynomial-size all-robust batches, and vanishing discarded-fraction tests for
arbitrary sampled batches under summability of \(\ell_m\nu_{d_m,a}^{(2)}\), with
the Bernstein replacement when the truncated mass is \(O(1/\ell_m)\).  The finite inverse-moment
form also carries the
random-search, success-count, Hamming-stability, robust-fraction, and
polynomial-batch robust-certificate package under the same matrix exception, and
allows arbitrary-size sampled batches with discarded fraction
\(\rho\ge e d^{-\tau}\).
Its schedule form
gives all-integer and polynomial
dimension consequences from polynomial inverse-moment decay, sparse
stretched-exponential consequences from logarithmic decay, and the critical
\(\log^3d\) integer-window consequence from \(O(1/\log d)\) truncated mass.
A power-curve
sufficient condition now makes this
checkable directly from one matrix: if
\(\Pr\{\|T_d(\xi)\|_{\mathrm{op}}<t aN_d\}\le L\ell^{-\gamma}t^\beta\) with
\(\beta>2\), then \(\mu_{d,a}^{\rm dy}\le C_\beta L\ell^{-\gamma}\); at
\(\gamma=1\), top-scale low-normalizer mass \(O(1/\log d)\) is compatible with
the Chernoff transfer whenever \(R>C_\beta L\).  The same power-curve exception
now carries the finite random-search, success-count, and Hamming-stable
robust-center package, so the weaker critical-density input yields actual
randomized certificates, not only a deterministic existence bound.  Its finite
abundance, robust-fraction, and polynomial-batch forms also match the Chernoff
branch: \(O_{\rho,S}(1)\) sampled signs certify a \(1-\rho\) robust fraction
with added \(d^{-S}\) failure, and \(M\le d^\beta\) batches are entirely robust
with added \(d^{-\sigma}\) failure under the same power-curve matrix exception.
Its
Borel--Cantelli form gives eventual randomized certificates whenever
\(\sum_m(d_m^{-A}+\ell_m^{1-\gamma})<\infty\); in the critical density
\(\gamma=1\), the Bernstein form instead needs only
\(\sum_m[d_m^{-A}+\exp\{-c n_m/\ell_m^2\}]<\infty\), giving all-integer
eventual certificates from \(n(d)\gtrsim\log^3 d\).  Adding
\(\sum_m d_m^{-\alpha}<\infty\) upgrades this to actual sampled-list success:
\(O_\alpha(1)\) random signs per dimension eventually contain robust
certificates almost surely, including in the critical power-curve regime.  The
same Borel branch now carries the finite abundance package: after adding
\(\sum_m d_m^{-S}<\infty\), \(O_{\rho,S}(1)\) sampled signs eventually have a
\(1-\rho\) robust-center fraction, and after adding
\(\sum_m d_m^{-\sigma}<\infty\), every sign in any batch \(M_m\le d_m^\beta\)
is eventually robust.  More generally, a power-curve sampled-count theorem allows
growing batches \(M_m\) and thresholds \(r_m\): if the sampled-list tail and
\(\binom{M_m}{r_m}d_m^{-B r_m}\) are summable, then eventually the whole list is
controlled and at least \(M_m-r_m+1\) sampled signs are good robust centers.  The
explicit full-integer critical corollary therefore
gives the same robust-fraction conclusion for \(S>1\), and all-robust polynomial
batches for \(\sigma>1\), from the same \(n(d)\gtrsim\log^3 d\) window; its
sampled-count variant also allows arbitrary batch sizes with a discarded
fraction \(\rho_d\to0\) whenever \(\rho_d\gtrsim d^{-\tau}\) for any
\(\tau<B-1\).  On stretched-exponential
dimension schedules \(d_m\ge\exp(c m^\kappa)\), this becomes explicit:
\(\gamma>1+1/\kappa\) is enough in the Markov form, while the critical
\(\gamma=1\) Bernstein form follows from any \(n_m\gtrsim\ell_m^{2+\rho}\),
\(\rho>0\); because \(\sum_m d_m^{-\alpha}<\infty\) for every \(\alpha>0\) on
such schedules, the same sparse windows give actual sampled-list robust
certificates, robust fractions, and all-robust polynomial batches from the
corresponding finite random sign trials.  The growing-count sparse variant also
handles polynomial-in-\(d_m\) and subcritical stretched-exponential sampled
batches after the binomial count tail is summable; taking
\(r_m=\lceil\rho M_m\rceil\) gives an eventual \(1-\rho\) robust-center fraction
for arbitrary sampled batch sizes, with only the list maximum retaining its
\(\log M_m\) entropy factor.  Quantitatively, the finite-size profile failure is
bounded by \(C d^{-A}+p_d(a)\min\{n,\log^2(e+d)\}\), and the same exceptional
probability now carries the finite random-search and logarithmic Hamming-stable
center package.  Thus a batch of \(M\) independent sign trials has the
entropy-sensitive maximum bound and amplifies the square-root failure to
\(d^{-AM}\), with matching binomial tails for successful certificates and robust
centers; equivalently, \(O_B(1)\) sampled signs find a logarithmically robust
certificate with added failure \(d^{-B}\).  The finite non-Chernoff package now
also has the same abundance, robust-fraction, and polynomial-batch forms as the
Chernoff branch: outside the same matrix exception and explicit sign tails,
\(O_{\rho,S}(1)\) sampled signs certify a \(1-\rho\) robust fraction with
additional \(d^{-S}\) failure, and every sign in a batch \(M\le d^\beta\) is
robust with additional \(d^{-\sigma}\) failure.  On sparse power-curve,
logarithmic-window, and Chernoff schedules, the same
\(r_m=\lceil\rho M_m\rceil\) count trick gives eventual \(1-\rho\)
robust-center fractions for arbitrary sampled batch sizes, again with only the
whole-list maximum retaining its \(\log M_m\) factor; in the sparse power-curve
and Chernoff count branches the discarded fraction may even vary with \(m\) and
tend to zero under any power threshold \(\rho_m\gtrsim d_m^{-\tau}\) with
\(\tau<B\).
Consequently a polylogarithmic
lower-tail rate \(p_d(a)\lesssim(\log d)^{-B}\), \(B>2\), gives a profile
failure \(O((\log d)^{2-B})\), while in \(n=O(\log d)\) windows \(B>1\)
suffices, and polynomial tails give polynomial failure with any smaller
exponent; these same explicit rates now apply to the full random-search and
robust-center package.  The summable version gives almost-sure eventual edge profiles along
triangular arrays, and carries the same random-search and logarithmic
Hamming-stable robust-center package eventually on the same matrix array; with
\(\sum_m d_m^{-\alpha}<\infty\), \(O_\alpha(1)\) sampled signs per dimension
themselves eventually find robust certificates almost surely over both the
matrices and sign samples.  The non-Chernoff count form also handles growing
sampled batches: once the sampled-list maximum and binomial count tails are
summable, all but the prescribed number of sampled signs are eventually good and
robust; under polynomial lower-normalizer tails, polynomial-size batches are
eventually entirely robust along the full integer sequence whenever the
sign-count exponent beats the batch growth by more than one, and along
polynomial dimension schedules \(d_m\gtrsim m^\kappa\) whenever
\((A-\beta)\kappa>1\).  Taking \(r_d=\lceil\rho M_d\rceil\) instead gives
eventual \(1-\rho\) robust-center fractions for arbitrary sampled batch sizes
along the full integer sequence as soon as the count exponent is \(A>1\), and
the same full-integer polynomial-tail count theorem lets \(\rho_d\to0\) whenever
\(\rho_d\gtrsim d^{-\tau}\) for any \(\tau<A-1\), again with no upper bound on
the sampled batch size.
Polynomial lower-normalizer tails \(p_d(a)\lesssim d^{-B}\)
with \(B>1\) need no dimension sparsification: the profile holds eventually
along the full sequence of integer dimensions.  More generally, polynomial tails
make this automatic along polynomial dimension schedules
\(d_m\gtrsim m^\kappa\) whenever \(B\kappa>1\), after choosing the auxiliary
matrix-event exponent \(A\kappa>1\).  The actual sampled-list version on the
same polynomial schedules adds only \(\alpha\kappa>1\), so \(O_\alpha(1)\)
random signs eventually contain robust certificates.  The matching
robust-fraction version needs only \(S\kappa>1\) for the fraction-confidence
summability, and under polynomial lower-normalizer tails it works in every
admissible sample window.  Polynomial lower-normalizer tails, as well
as logarithmic tails with exponent \(B>3\), make the same conclusion automatic
along exponentially sparse dimension schedules; because \(\sum_m d_m^{-\alpha}\)
is then summable for every \(\alpha>0\), the sparse exponential schedules also
give actual sampled-list robust certificates from a constant number of sign
trials.  In logarithmic sample windows
\(n_m\lesssim\log^\gamma d_m\) along stretched-exponential schedules
\(d_m\ge\exp(c m^\kappa)\), the sparse almost-sure logarithmic threshold
improves to \(B>\min\{\gamma,2\}+1/\kappa\); on log-corrected schedules
\(d_m\ge\exp(c m^\kappa(\log m)^\lambda)\), the endpoint
\(\kappa(B-\min\{\gamma,2\})=1\) is allowed when
\(\lambda(B-\min\{\gamma,2\})>1\); if this endpoint is itself critical, a
second \(\log\log m\) correction recovers the next integral-test boundary.  A
summability-only sampled-count form covers arbitrary schedules with
\(\sum_m(\log(e+d_m))^{\min\{\gamma,2\}-B}<\infty\), equivalently it is enough
to design a lower logarithmic profile \(h_m\le\log(e+d_m)\) with
\(\sum_mh_m^{\min\{\gamma,2\}-B}<\infty\).  In the raw summability form, fixed
robust fractions for arbitrary sampled batches need only
\(\sum_m d_m^{-(A-\tau_0)}<\infty\) for some \(0<\tau_0<A\), in addition to the
logarithmic matrix and sampled-list tails.  The same schedules give actual
sampled-list robust certificates from \(O_\alpha(1)\) signs for every
\(\alpha>0\), vanishing discarded-fraction robust-center tests for arbitrary
sampled batches, and the log-corrected sampled-count forms support all-robust
polynomial-in-\(d_m\), subcritical log-corrected, and second-order log-corrected
stretched-exponential batches.  The matching
robust-fraction version gives, for every \(S>0\)
and \(0<\rho\le1\), a constant-size sampled list with an eventual \(1-\rho\)
fraction of logarithmically robust centers.  The matching sparse count version
also supports growing batches:
polynomial-in-\(d_m\) batches are eventually entirely robust whenever
\(\beta<A\), and \(M_m\le\exp(bm^\sigma)\) is allowed when
\(\sigma<\kappa\), or \(\sigma=\kappa\) and \(b<Ac\).
In the large-sample window
\(n/\log^2 d\to\infty\), there is a
tunable Chernoff tradeoff: if \(p_d(a)\le\theta/\log^2 d\) and the empirical
budget is \(Rn/\log^2 d\) with \(R\ge1\) and \(R>\theta\), the excess-count cost is
\(\exp\{-I(R,\theta)n/\log^2 d\}\), \(I(R,\theta)=R\log(R/\theta)-R+\theta\).
In particular \(p_d(a)\le L/\log^2 d\) permits any relative excess budget
\((1+\eta)Ln/\log^2 d\) with cost
\(\exp\{-Lh(\eta)n/\log^2 d\}\), \(h(\eta)=(1+\eta)\log(1+\eta)-\eta\);
taking \(\eta=1\) gives profile failure
\(O(d^{-A}+\exp\{-c_Ln/\log^2 d\})\), and the same Chernoff-improved
exception controls the random-search and robust-center certificate package.
The finite Chernoff search package is now joint in matrices and signs:
sampled-list maxima add only the \(e^{-u}\) sign tail, while good-sign and
robust-center counts add the same binomial tails as in the deterministic
edge-profile theorem.  In particular, \(O_B(1)\) sampled signs find a robust
certificate with additional failure \(d^{-B}\); more generally, for any batch
and fraction \(\rho\), the sampled-list maximum and a \(1-\rho\) robust fraction
hold simultaneously outside the explicit binomial tail
\(\binom{M}{\lceil\rho M\rceil}d^{-B\lceil\rho M\rceil}\).  Choosing
\(M=O_{\rho,S}(1)\) and the count exponent large enough gives the same
\(1-\rho\) robust fraction with additional failure \(d^{-S}\), and any finite batch
\(M\le d^\beta\) is entirely robust with additional failure \(d^{-\sigma}\)
after choosing the count exponent above \(\beta\).  The corresponding Borel count
form handles growing batches: under the summable maximum and binomial
failure-count tails, all but the prescribed number of sampled signs are
eventually robust centers.  In particular the constant-size
\(1-\rho\)-robust-fraction test persists almost surely whenever
\(\sum_m d_m^{-S}<\infty\); on polynomial dimension schedules
\(d_m\gtrsim m^\kappa\), this is the condition \(S\kappa>1\).  Polynomial-size batches are eventually entirely
robust when the sign-tail exponent beats the batch growth by more than one.  The
summable Chernoff version gives almost-sure eventual edge profiles, signing
certificates, scalar moments, random search, and robust-center certificates under
\(\sum_m(d_m^{-A}+\exp\{-c n_m/\log^2 d_m\})<\infty\).  Adding
\(\sum_m d_m^{-\alpha}<\infty\) gives actual sampled-list robust certificates
from \(O_\alpha(1)\) random signs per dimension; along all integer dimensions
this adds the condition \(\alpha>1\).  Equivalently, once
\(\sum_m d_m^{-A}<\infty\), the critical-log window
\(n_m\gtrsim\log^2(d_m)\log m\) is enough, and on the full integer sequence
\(n(d)\gtrsim\log^3d\) also supports arbitrary sampled batches with robust-center
fractions tending to one when \(\rho_d\gtrsim d^{-\tau}\) for any \(\tau<B-1\);
along stretched-exponential
dimension schedules any window \(n_m\gtrsim\log^{2+\rho}d_m\), \(\rho>0\), is
enough, and the same sparse schedules give actual sampled-list certificates for
every \(\alpha>0\).  They also give the constant-size \(1-\rho\)
robust-fraction test for every confidence exponent \(S>0\), removing the
all-integer summability restriction \(S>1\).  The sparse count version also improves growing-batch
search: if \(M_m\le d_m^\beta\), then every sampled sign is eventually robust
as soon as the sign-count exponent satisfies \(B>\beta\), and batches
\(M_m\le\exp(bm^\sigma)\) are allowed whenever \(\sigma<\kappa\), or
\(\sigma=\kappa\) and \(b<Bc_0\).  Along all integer dimensions
\(n(d)\gtrsim\log^3 d\) suffices, with \(\alpha>1\) for the actual sampled-list version.
Convex
concentration verifies this condition, so the
full realization-normalized theorem holds for real or complex Gaussian and
bounded coefficient laws, for Rademacher coefficients, and now also for
strongly log-concave real or complex coefficient laws, when
\(n\le d\le e^{o(n)}\) and \(d/n\to\infty\). A finite-convolution closure also
covers fixed independent linear combinations of convex-concentrated laws, such
as Gaussian plus Rademacher or bounded noise, after variance normalization.
Bounded Gaussian scale mixtures with volatility bounded above and below are
covered by a separate conditional-Gaussian lower-tail argument. Arbitrary
positive matrix-level radial amplitudes are also allowed, with no moment or
tail assumptions, because they cancel pathwise under realization normalization.
The heterogeneous lower-normalizer library theorem also removes the identical-law
restriction: the matrices may be drawn from different laws in any uniform
convex, strongly log-concave, finite-convolution, bounded Gaussian
scale-mixture, or other polynomial lower-normalizer library, and finite
matrix-level mixtures are included by conditioning on the selected laws. The
uniform-library input can now be relaxed
to an aggregate lower-tail budget
\(\sum_i p_{d,i}(a)\min\{1,\log^2(e+d)/n\}=o(1)\), allowing sparse bad-law
contamination or nonuniform law quality. A sharper heterogeneous dyadic-budget
variant charges each bad law by depth below the cutoff and only requires
\(\log(e+d)n^{-1}\sum_i\mathbb E Y_{d,i,a}^{\rm dy}\to0\), allowing top-scale
low-normalizer contamination at density \(o(1/\log d)\). The dyadic budget is
checkable from the aggregate truncated inverse-square normalizer mass
\(\log(e+d)n^{-1}\sum_i\nu_{d,i,a}^{(2)}\to0\), with the same explicit
inverse-\(p\)-moment loss as in the iid theory. Its Borel sampled-operator
schedule gives eventual Toeplitz/Hankel robust centers under summable aggregate
inverse-moment matrix failure and sampled-count tails; the extracted fraction
form gives fixed robust fractions for arbitrary sampled batch sizes, vanishing
discarded fractions, and all-robust polynomial batches under the same matrix
budget plus a summable sign-tail power. A batch-maximal heterogeneous
inverse-moment form additionally keeps the whole sampled-list maximum at the
entropy-sensitive scale, both on the edge grid and in full Toeplitz/Hankel
operator norm; its schedule corollary gives explicit all-integer,
polynomial-dimension, stretched-exponential, and Bernstein-window thresholds for
fixed fractions, vanishing discarded fractions, and all-robust polynomial
batches. The inverse-\(p\) schedule version replaces the aggregate inverse-square
budget by explicit average inverse-\(p\) decay, with stretched-exponential
threshold \(B>1+(2-p)_+/2+1/\kappa\). A random-environment version permits the
coefficient laws themselves to be selected randomly, including random sparse
contamination patterns, provided the empirical inverse-\(p\) budget exceeds the
deterministic schedule only on summably many dimensions. The expectation
criterion verifies those tails from \(\sum m^B\beta_m\),
\(\sum d_m^B\beta_m\), or \(\sum\ell_m^B\beta_m\) summability, and countable
random law catalogs reduce to marginal catalog-weight averages. If the catalog
choices are independent, or are exposed predictably with deterministic one-step
envelopes, and the selected inverse-\(p\) weights are bounded by \(L_m\),
Bernstein/Freedman criteria replace the Markov loss by the sharper tail
\(\exp(-c n_m\alpha_m/L_m)\); a truncated Freedman variant also covers
unbounded predictable catalogs after adding a summable excess-tail budget, with
conditional moment envelopes giving a direct check of that excess and an
automatic truncation choice reducing it to one moment series. The
active-menu version makes the same automatic truncation simultaneous over
finite predictable-catalog menus, paying \(J_m\) in the environment,
matrix, and sampled-sign summability terms; the active-prefix schedule
corollary gives explicit all-integer, polynomial-dimension,
stretched-exponential, and logarithmic-prefix thresholds, and a random
active-count version adds only a summable overflow-tail requirement, with
sub-Poisson proposal-count envelopes and predictable one-step intensity budgets
providing direct overflow tests, even when the intensity budget itself is
certified outside a summable random exceptional set; moment tests make those
random intensity-envelope failures directly checkable, while bounded and
truncated scalar Freedman tests replace moments by predictable mean,
square-function, jump, and excess envelopes; automatic moment truncation then
gives concrete all-integer, polynomial-dimension, stretched-exponential, and
logarithmic-prefix rate tests for heavy-tailed predictable loads, and the
deterministic active prefix can be chosen directly as
\(\lceil\max\{2\nu_m,Q\log(e+m)\}\rceil\); an end-to-end
corollary couples this automatic prefix with the active-menu catalog moment
budget, so the sampled Toeplitz/Hankel certificates follow from the activation
and catalog moment series plus the matrix/sign tail schedules, with concrete
rate tests for all-integer, polynomial-dimension, stretched-exponential
polynomial-active, and logarithmic-active schedules, plus sample-size lower
bounds and unbounded-aspect compatibility windows under real aggregate
catalog-moment growth exponents, with per-stream row-catalog budgets converted
by the automatic prefix size and the remaining mean-budget parameter tuned from
the per-stream first-moment decay exponent.
The same
heterogeneous,
average-budget, dyadic-budget, and inverse-moment-budget inputs now have the
sampled operator-certificate package:
constant-size searches, fixed-size robust-fraction tests, polynomial-size
all-robust sampled batches, and almost-sure sampled-batch schedules. In these
same concentration classes,
constant-size sampled sign lists also contain
logarithmic-radius Hamming balls that are robust certificates for both full
Toeplitz and full Hankel operator norm, with probability tending to one; fixed
lists have \(1-\rho\) robust-center fractions, and polynomial-size sampled
batches are entirely robust with high probability. The same lower-tail input now
has an almost-sure sampled-batch schedule form: on summable schedules,
and in particular on stretched-exponential dimension schedules, polynomial-size
sampled batches are eventually entirely Toeplitz/Hankel operator-robust. Thus the
remaining open full-bandwidth problem is narrowed to arbitrary sub-Gaussian laws
without this aggregate lower-tail, dyadic-depth, or inverse-moment input, even
after allowing heterogeneous uniform libraries and sparse bad-law budgets, where
deterministic contraction tools still certify only
\(\sqrt{n\log(e+d/n)}\).  The paper
explains why the intrinsic-freeness theory of
Bandeira--Boedihardjo--van Handel does not by itself reach $O(\sqrt n)$
(Toeplitz matrices are Hilbert--Schmidt delocalized, so the free-deviation
correction is not lower order), and adds an edge-variance criterion: if the
realization-normalized Fourier samples have empirical square mass
\(O(n/\log d)\) at every edge node, then independent signs already give
\(O(\sqrt n)\).  The residual input is therefore a one-matrix low-normalizer
tail beyond these concentration, finite-convolution, and bounded scale-mixture
classes, not another generic discrepancy
theorem.  The moment section now proves exact fixed-edge fourth- and
sixth-moment identities.  The fourth identity is
\[
\mathbb E_\varepsilon\Big|\sum_i\varepsilon_i a_i\Big|^4
=2\Big(\sum_i|a_i|^2\Big)^2+\Big|\sum_i a_i^2\Big|^2
-2\sum_i|a_i|^4,
\]
and proves that in the expected-normalized model the phase-square term
\(\max_{\omega^d=1,\omega^2\ne1}|\sum_i v_i(\omega)^2|\) is
\(O(\sqrt{n/\log d})\) with high probability, yielding the uniform
expected-normalized fourth-moment profile
\(\max_{\omega^d=1}\mathbb E_\varepsilon|\sum_i\varepsilon_i v_i(\omega)|^4
 =O((n/\log d)^2)\) and the non-real complex improvement with leading constant
2.  The sixth identity adds only the weighted phase-square statistic and
diagonal fourth/sixth masses, gives a perturbative criterion for the complex
constant \(6\), and combines with phase-square cancellation plus a uniform
expected-normalized root-entry bound to give
\(\max_{\omega^d=1}\mathbb E_\varepsilon|\sum_i\varepsilon_i v_i(\omega)|^6
=O((n/\log d)^3)\) in the expected-normalized model, with the non-real
refinement
\(\mathbb E_\varepsilon|\sum_i\varepsilon_i v_i(\omega)|^6
=6(\sum_i|v_i(\omega)|^2)^3+O((n/\log d)^2)\).  A final scalar moment corollary
gives, simultaneously for every \(k\ge1\),
\[
\max_{\omega^d=1}\mathbb E_\varepsilon
  \Big|\sum_i\varepsilon_i v_i(\omega)\Big|^{2k}
\le \big(Ck\,n/\log d\big)^k
\]
on the same expected-normalized edge-profile event.  A deterministic max-edge
moment corollary sums these scalar bounds over the \(d\) edge nodes, choosing
\(k\asymp\log d\) recovers the high-probability edge-control criterion by Markov,
and the expected-normalized edge-maximum package gives an \(O(\sqrt n)\)
sign-averaged edge maximum, \(d^{-B}\) sign-tail, deterministic signing, and
constant-size sampled robust certificates with logarithmic Hamming stability on
the same high-probability matrix event.  Thus the residual higher-moment route must control the higher
pairing analogues uniformly over all edge nodes in the realization-normalized
model, not merely scalar edge moments under the variance profile.  The meeting
remark separates the
disk-count fluctuation problem of paper 5 from the one-symbol Gaussian boundary
small-ball input proved in paper 16.

\textbf{15. \texttt{UDC\_General\_Position\_v1.pdf}} (13pp) and
\texttt{UDC\_Curved\_Truncations\_v1.pdf} (19pp, Thread C).
The interpolation law beyond the orthogonal model for general-position flat
truncations, with a fixed-family curved companion and curvature-band-uniform
curved lower/slab constants left open.  The first paper treats a $C^2$ strictly convex body with curvatures
in a fixed band, truncated by $k$ half-spaces in general position along a
$j$-stratum ($j=n-1-k$), and proves the critical scaling is
$\Theta((\log N/N)^{2/(2n-j)})$, exponent and logarithm exact; the mechanism
replacing the orthogonal model's axis windows and reflection fraction $2^{-k}$
is the conjugate frame (dual basis to the truncating normals) and the normalized
solid angle of the normal cone, with a proof that John-position rounding
provably cannot reduce to the orthogonal case.  The companion removes the
remaining flatness by treating a \emph{fixed} smooth strictly convex truncation
family: the stratum cap still has volume $\asymp h^{(2n-j)/2}$ and the
constraint slice stays admissible with fixed-configuration constants (closing the slicing induction), via a
steep/shallow trichotomy at the angular threshold $\sqrt h$, so the same law
holds for fixed curved data.  A constant uniform over all curved families with
curvatures in a fixed band, near the multiply-tangent directions, is stated as
the residual open estimate.  The curved paper now also proves a conditional
uniformity theorem under resolved multi-tangent geometry (uniform slice-curvature
floor plus projected-normal Gram lower bound), records that every finite atlas
and every compact resolved atlas of fixed curved configurations is uniform with
atlas-dependent constants, proves that resolved geometry is stable under
\(C^2\)-small perturbations, so every compact \(C^2\)-continuous resolved atlas
is uniform by a finite-subcover argument, extends this to compact-stage
uniformity on \(\sigma\)-compact exhaustive resolved atlases, and
proves a fully uniform one-curved-cut corollary when the cut splays quantitatively away
from \(\partial K\), now
extended to the curvature-dominated oblique one-cut regime
\(c_K-\max(\langle\nu_K,n_1\rangle,0)C_g/\gamma_0>0\), with
nonobtuse contact \(\langle\nu_K,n_1\rangle\le0\) as an explicit corollary.
The remaining
uniformity question is localized to one-cut curvature-separation failures and
genuine multi-tangent degeneration.

\textbf{16. \texttt{Banded\_Toeplitz\_Conditioning\_v1.pdf}} (102pp, Thread B).
The conditioning (least-singular-value) question for random balanced banded
Toeplitz matrices, addressing the balanced-regime gap left open by paper 5.  The
current draft proves the finite-section reduction on the zero-winding/stable
event, a sharp Kac--Rice Gaussian single-arc boundary small-ball estimate away
from \(z=\pm1\), the matching away-from-endpoint first moment for crossings, an
explicit fixed-bulk Kac--Rice down-crossing density
\(t\sqrt N/\sqrt{12\pi}\) per unit arclength, and the corresponding
continuum-scale fixed-bulk limit
\((2\pi-4\delta)\tau/\sqrt{12\pi}\) at \(t=\tau/\sqrt N\),
the resulting fixed-bulk first-moment probability upper bound with the same
down-crossing constant, and a Paley--Zygmund transfer showing that a
second-factorial-moment bound for the fixed-bulk down-crossing count gives a
matching small-\(\tau\) fixed-bulk lower tail with leading constant
\((2\pi-4\delta)/\sqrt{12\pi}\). Combining the two gives the fixed-bulk
first-order probability law as \(\tau\downarrow0\); combining the fixed-bulk
upper bound with endpoint-neighborhood separation gives the sharp whole-circle
continuum upper constant \(\sqrt{\pi/3}\,\tau\), and if the same factorial
moment bound is uniform as \(\delta\downarrow0\), the whole-circle first-order
probability law follows as \(\tau\downarrow0\). The sharp global upper tail also
implies that the inverse scaled minimum
\((\sqrt N\min_{\mathbb T}|G_m|)^{-1}\) is uniformly weak \(L^1\), with
finite-window sharp weak constant \(\sqrt{\pi/3}\), gives inverse upper-tail
quantiles \(O(\alpha^{-1})\) with sharp limsup constant
\(\sqrt{\pi/3}\,\alpha^{-1}\), gives simultaneous finite-batch weak-\(L^1\)
and quantile bounds with the same sharp constant times the batch size, gives
finite-batch bad-clearance count and order-statistic quantile bounds, summable
first-bad, total bad-mass, jump/Stieltjes-weighted and powered clearance-deficit mass,
normalized deficit-budget, subcritical inverse-clearance mass, Stieltjes inverse-gauge mass,
confidence/menu and aggregate-menu budgets, and sparse bad-count growing-batch clearance schedules,
subcritical inverse-power bounds, and critical
truncated/logarithmic inverse-moment bounds with the same linear batch-size
factor, gives critical
truncated inverse expectation \(O(\log R)\) with the
same sharp finite-window coefficient, gives uniform
logarithmic inverse moments of every positive order with limsup constant at
most \(\sqrt{\pi/3}\Gamma(p+1)\), and hence uniform negative
moments of every order \(q<1\) for
\(\sqrt N\min_{\mathbb T}|G_m|\), with small-window coefficient
\(\sqrt{\pi/3}/(1-q)\). It also includes the
endpoint Kac--Rice estimate, endpoint lower sharpness, the exact bivariate
two-endpoint Gaussian formula and sharp second-order expansion
\(\mathbb P\{\min(|G_m(1)|,|G_m(-1)|)\le t\}
=2\sqrt{2/\pi}\,t/\sqrt N-(2/\pi)t^2/\sqrt{N^2-\rho_N^2}
+O(t^3/N^{3/2})\), where \(N=m+1\) and
\(\rho_N=\sum_{j=0}^{N-1}(-1)^j\).  At \(t=\tau/\sqrt N\) this gives endpoint
probability \(2\sqrt{2/\pi}\,\tau/N+O_\tau(N^{-2})\), separating the parity
endpoints from the \(O(1)\) non-real continuum law; the full natural
\(O(1/m)\) endpoint neighborhoods also contribute only \(O_{C_0}(\tau/N)\).
The draft also proves the
global endpoint floor, the sharp uniform
Gaussian boundary small-ball law, bounded-density fixed-point,
deterministic-grid, deterministic and independent-random subquadratic-grid-miss,
quadratic, independent uniform random-grid, stratified random-grid, variance-only, and
speed-tail grid detection,
and speed-conditional local/uniform
anti-concentration, the
exact Gaussian \(N\)-th-root-grid law, including the full-grid endpoint
survival factors, showing the \(t^2\) pointwise grid artifact and its
fixed pointwise Rayleigh moment expansion and exact \(N^{-1}\) continuum-scale grid
coefficient, the
reduction from a polynomial speed tail to the polynomial-scale boundary law,
all-polynomial-moment and subexponential bounded-density verifications of that
speed-tail input, the variance-only polynomial derivative cutoff, tail-free
polynomial boundary separation, fixed-power and arbitrary-confidence
polynomial-grid detection, plus the resulting sampling bracket, the
bounded-density conditioning theorem under circle clearance, and the
endpoint/order-\(q\) root-of-unity, fixed-grid, and bounded-denominator
Rademacher lattice anti-concentration scales, the matching \(N^{-1}\)
fixed-radius lower bound at every fixed non-real rational angle, and the exact
fixed-root atom law: at a primitive \(q\)-th root the zero atom has a local
limit \(\gamma_qN^{-\varphi(q)/2}\), with an explicit kernel-covolume constant,
on the admissible parity residue classes and is absent on the inadmissible
classes; finite sets of cyclotomic orders have leading exact-zero asymptotics
on their first nonempty totient layer; consequently fixed bounded-denominator
non-real grids have a peeled exact-zero coefficient after removing all primitive
orders below any fixed cyclotomic degree threshold.  The unpeeled fixed-\(Q\)
union is now expanded through the second order:
\(c_Q(N)/N+e_Q(N)/N^2+o_Q(N^{-2})\).  The leading coefficient \(c_Q(N)\) is the
sum of the primitive quadratic pair constants: \(2\sqrt3/\pi\) from each of
\(q=3,6\) on \(3\mid N\), and \(4/\pi\) from \(q=4\) on \(4\mid N\).  The
second coefficient \(e_Q(N)\) combines the second-order corrections to those
three single atoms, their pairwise overlaps
\((48/\pi^2\) for \(E_3\cap E_6\) on \(6\mid N\), and \(8\sqrt3/\pi^2\) for
\(E_3\cap E_4\) and \(E_4\cap E_6\) on \(12\mid N)\), and the first
nonquadratic layer.  After removing cyclotomic degree two, that layer is also
explicit: primitive orders \(5,8,10,12\) give the complete \(N^{-2}\)
coefficient on their admissible divisibility classes.  Adding the two endpoints
gives the full bounded-denominator root-of-unity exact-zero expansion through
order \(N^{-1}\): an exact central-binomial endpoint-union term plus the same
non-real coefficient \(c_Q(N)/N\).
Therefore, for standard Gaussian coefficients, the arbitrary fixed
polynomial-in-\(n\) conditioning conclusion on the event \(\{\nu=r\}\) is now
unconditional, using the Konyagin--Schlag circle-clearance estimate plus a
deterministic uniform finite-section threshold under a clearance floor.  For
non-Gaussian bounded-density laws satisfying the same
circle-clearance input, the same finite-section conclusion now follows without
any tail assumption beyond variance one; the finite-grid, speed-cutoff
continuum, variance-only boundary separation with explicit
\(O_{A,K}(m^{-2(A-1)/3})\) exceptional probability and grid detection with
\(O_\delta(m^{1-2\delta})\) miss probability, the finite-confidence
variance-only derivative cutoff
\(\Pr\{\sup_{\mathbb T}|G_m'|>R\}\le Cm^4/R^2\), its
\(\eta\)-level grid detector at mesh
\(O(\tau\sqrt\eta/(m^2\sqrt N))\), and its bounded-density boundary loss
\(\eta+C_{B,K}mt/\sqrt\eta\), fixed-power
\(O_\beta(N^{5-2\beta})\) miss rates for \(\beta>5/2\), equivalently
\(O_A(m^{-A})\) after \(O_{A,\tau}(N^{(A+5)/2})\) grid points, and under a
natural speed tail the sharper \(O_{A,\tau}(N^2\sqrt{\log N})\) grid
sufficiency with the same polynomial miss probability, while independent
uniform random grids have the corresponding
\(O_{A,\tau}(N^2(\log N)^{3/2})\) sufficiency and stratified one-point-per-cell
random grids avoid that coupon overhead at the deterministic
\(O_{A,\tau}(N^2\sqrt{\log N})\) scale; speed-tail
boundary/grid reductions, deterministic sampling brackets, all-polynomial-moment speed-tail, subexponential
speed-tail, variance-only cutoff, and
bounded-density-conditioning theorem under clearance are proved.  The stable-event
conditioning loss is now quantitative and tunable: if the clearance event fails
with probability \(q_m\), the additive loss is
\(q_m+C_{A,K}m^{-2(A-1)/3}+(m+1)/R_m^2\), hence the coefficient-size term may be
made \(O(m^{-D})\) for any fixed \(D\) after enlarging only the deterministic
threshold; choosing the boundary floor and envelope together gives the packaged
loss \(q_m+O_D(m^{-D})\), and conditioning on \(\{\nu=r\}\) divides this loss by
\(\Pr\{\nu=r\}\).  The
cutoff-free non-Gaussian speed/crossing law remains open for the full
polynomial-scale boundary small-ball theorem.  Sharp
constants, anti-clustering/decorrelation for the global two-sided boundary
minimum law, continuum non-Gaussian universality at full bounded-density
generality, and the continuum/local-CLT discrete analogue beyond
bounded-denominator lattice scales remain open, and the result is not an
unconditional high-probability resolution of balanced stability.

\textbf{17. \texttt{Master\_File\_for\_Solving\_IV\_Surface\_and\_Options\_Open\_Problems.pdf}}
(286pp, Thread D).  A new options-market master file saved in
\texttt{paper17/}.  It now includes a repo-local imported theorem bank under
\texttt{paper17/imported\_results/}, with a manifest distinguishing verified
packets from candidates still needing a final skeptic gate.  It records the
local IV solved-result bank from \texttt{iv\_stoch\_research} and companion
workspaces (SSVI/eSSVI no-arbitrage, Lee/Benaim--Friz wings,
BBF short-time harmonic mean, one-half slope, Heston ATM skew, local-variance
Jensen gap, variance/VIX convexity adjustments, CIR spectral skew decay,
CIR-versus-lognormal-vol discrimination, rough-volatility short-time skew, and
large-maturity flattening), states 100 stable open problems across static
surface geometry, IV dynamics, short-maturity/event asymptotics, rough/Volterra
volatility, wings and moment explosions, IV percentiles, gamma and hedging,
robust/MOT pricing, quote microstructure, and learning/calibration/formal
verification, proves the finite-grid arbitrage-free projection theorem for
quote repair, adds a distribution-free finite-sample theorem package for
surface IV percentile ranks and confidence bands, adds a finite-mesh
gamma-carry theorem package with transaction-cost lower bounds and a local-time
limit, adds a finite-scenario robust-pricing/MOT package with bid-ask surface
uncertainty, frictional hedge libraries, roughness budgets, and quote-error
multiplier allocation, adds a distribution-free opening-price and opening-IV
surface package using conformal prediction bands, Lipschitz price transfer, and
overnight tail reserves, adds a finite bid-ask no-arbitrage feasibility and
liquidity-weighted repair package with Farkas certificates, interval relaxation,
and one-quote influence bounds, adds a finite-sample density-recovery package
with no-arbitrage confidence intervals for strike-bin masses and a lowest-strike
atom bound, adds a certified implied-volatility inversion package with
price-interval inversion, vega-stability, and Greek interval propagation, adds a
finite-library data-snooping control package with holdout lower confidence
bounds, block validation, and familywise signal gates, adds a latency-aware
surface risk package with deterministic and high-probability stale-surface
pricing reserves, Greek reserves, and execution gates, adds an exchange-comparable
surface-coordinate package with forward-normalized call coordinates, gauge
invariance, common venue polytopes, and repair-stability bounds, adds a short-end
skew-exponent diagnostics package with log-slope confidence sets, Hurst intervals,
second-order extrapolation bands, and finite model-class rejection certificates,
adds a finite wing-envelope and tail-reserve package with no-arbitrage wing
completion polytopes, certified finite-strike wing-slope intervals, cutoff
variance-swap reserves, and tail-robust calibration gates, adds a robust VIX
convexity-envelope package with bucketed moment/tail LP bounds for
\(\mathbb E[\sqrt W]\), finite-support exact bounds, convexity-gap reserves, and VIX
calibration gates, adds a monitored-barrier and IV-factor superhedging package
with finite IV-factor path-market duality, monitored barrier price intervals,
continuity-bracket reserves, and marginal value tests for tradable IV factors,
adds a robust Bermudan/American exercise package with finite stopping-policy
reduction, robust seller and buyer exercise intervals, exercise-policy
certificates, and mesh-to-continuous exercise reserves, adds a hard-shape
learning-layer package with no-arbitrage projection layers, convex-template
layers, approximation inheritance, and finite-sample risk bounds, adds an
uncertainty-aware surface-ensemble package with convex no-arbitrage ensembles,
conformal surface balls, Lipschitz price intervals, and projection-stable
preprocessing, adds a theorem-fingerprint model-selection package with
finite fingerprint boxes, convex rejection certificates, stable relaxed gates,
and survivor holdout selection bounds, adds a differentiable-pricing error
certificate package with implemented-graph AD, deterministic truncation
intervals, Monte Carlo price/Greek confidence radii, and book-level aggregation,
adds an end-to-end theorem-to-trading audit package with finite certificate
ledgers, robust interval action gates, realized P\&L residual intervals, and
replayable pipeline certificates, adds a finite Wasserstein calibration package
with transport-loss existence, empirical stability, identifiable-parameter
consistency, misspecification bounds, and calibration-to-price intervals, adds a
proof-carrying finite no-arbitrage verification package with exact rational
acceptance/rejection certificates, finite Breeden--Litzenberger checks, and
deterministic replay ledgers, adds a finite no-arbitrage neural-operator package
with projected approximation, parameter-covering generalization, empirical-risk
oracle bounds, and deployment certificates, adds a finite calibration-landscape
topology package with compact sublevel stability, identifiable-basin
localization, finite graph bracketing, and local-minimum exclusion certificates,
adds a finite eSSVI projection-certificate package with compact-image
existence, distance stability, finite-net global certificates, and conditional
uniqueness under quadratic separation, adds a finite SSVI density-frontier
package with residual representation, mesh-certified safe/unsafe bands,
parameter-distance margins, and cellwise frontier localization, adds a finite-to-continuum
no-arbitrage transfer package with common-grid interpolation, node-noise margin
conditions, projection-to-continuum repair, and explicit piecewise-linear
density/calendar measures, adds a finite missing-tenor interpolation package
with exact insertion polytopes, LP value bands, projection repair,
infeasibility certificates, and node-noise acceptance margins, adds a finite
no-arbitrage polytope regularity package with active tangent and normal cones,
sharp feasible step lengths, KKT calibration certificates, active-face
stratification, and finite zero-curvature geometry, adds a finite no-arbitrage
stress-deformation package with factor stress cones, exact amplitudes, safe
factor boxes, LP book-impact bands, projection repair, and smooth-factor
first-order transfer, adds a finite minimal static-statistics package with
facet-test minimality, redundancy LP/Farkas certificates, residual
distance-to-repair bounds, and exact finite/affine parametric-family gates,
adds a finite multi-asset consistency package with joint-state component/index
feasibility, Farkas separating portfolios, projection repair, and
dispersion/correlation LP bands, adds a finite dynamic surface no-arbitrage
package with scenario-tree martingale feasibility, drift-hull restrictions,
separating self-financing certificates, and diffusion-skeleton gates, adds a
finite minimal Markov surface-state package with calibrated-tree quotients,
feature-state exactness gates, cardinality lower bounds, and approximate
lumpability price bands, adds a finite recalibration-limit package with
projected no-arbitrage recursions, normal-cone repair accounting, stopped
interior diffusion limits, and linear covariance bands, adds a finite surface
covariance-operator package with active-facet admissibility, tangent covariance
projection, no-arbitrage PCA, shrinkage, and portfolio variance bands, adds a
finite forward-start smile package with two-date martingale-coupling
consistency, conditional forward-smile kernels, projection repair, and
cliquet/forward-volatility LP envelopes, adds a finite self-financing surface
drift package with discounted-price drift-span restrictions, factor-map Itô
drift gates, instantaneous separators, boundary-compatible drift tests, and
projection repair, adds a finite IV surface filtering package with hidden-state
no-arbitrage Bayes recursions, asynchronous quote likelihoods, posterior
no-arbitrage means, credible price bands, quote infeasibility certificates, and
robust action gates, adds a finite surface-jump structure package with static
no-arbitrage jump cones, event martingale hulls, event payoff LP envelopes,
jump-library amplitude polytopes, and repair multipliers, adds a finite surface
transport metric package with common-grid law extraction, maturity transport LPs,
no-arbitrage interpolation paths, transport shock-ball reserves, barycenter
regime prototypes, and repair-stability budgets, adds a finite no-arbitrage PCA
package with active-face factor admissibility, exact lineality-constrained PCA,
score polytopes, box-safe score radii, factor hedge residual bands, and raw-factor
repair/rejection certificates, adds a finite uniform BBF convergence package with
strike-window grid-to-window bounds, quote-error confidence bands, rate
certificates, adaptive mesh gates, and zero-DTE extrapolation intervals, and
adds a finite second-order one-half-slope package with certified ATM skew extraction,
two-term coefficient polytopes, zero-DTE skew extrapolation bands, and candidate
local-vol coefficient rejection, adds a finite short-time curvature package with
ATM-curvature extraction, curvature coefficient polytopes, zero-DTE curvature
bands, and butterfly/broken-wing price reserves, adds a finite event-time smile-jump
package with martingale event-kernel convolution, feasible kernel polytopes,
event variance/straddle bands, smile-jump bands, and event-calendar reserves, adds
a finite open-auction option-asymptotics package with indication-to-print kernels,
opening option/book price bands, opening-gap scaling certificates, and first-print
IV transfer, adds a finite barrier short-time correction package with terminal-law
mass intervals, local-vol bridge hit bands, barrier-hit LP certificates, and
knock-in/knock-out correction intervals, adds a finite zero-DTE gamma/cost package
with gamma-theta-cost scaling tubes, hedge-cadence cost bands, and
convexity-survival gates, adds a finite stochastic-rate short-smile package with
stochastic-carry cell LPs, rate/dividend price corrections, IV transfer bands, and
first-two-smile-term inclusion gates, adds a finite expiration-pinning package with
pinning-window mass LPs, point/node atom certificates, call-convexity spike tests, IV curvature
transfer, and smooth-density rejection gates, adds a finite rough-Heston skew-constant package with
signed constant tubes, pairwise maturity consistency tests, candidate formula audits, and
parameter-cell rejection certificates, adds a finite Volterra kernel-identifiability package with
kernel response polytopes, functional bands, invisible-nullspace certificates, and kernel-cell
rejection gates, adds a finite sum-of-exponentials error-propagation package with kernel-error
ledgers, state perturbation envelopes, price/Greek/IV reserves, and calibration-safe acceptance
gates, adds a finite rough-vol dynamic no-arbitrage package with variance-forward cash-mark
martingale feasibility, rolling-state bridge constraints, rough-kernel drift gates, bounded
stock-martingale certificates, and parameter-cell rejection tests, adds a finite second-order
rough-smile certificate package with coefficient intervals, reserve regression, Lean-ready
Hurst-cell envelopes, and zero-DTE IV/price gates, adds a finite rough-leverage identifiability
package with product tubes, signed-product ambiguity certificates, product-plus-magnitude lifted
feasible sets, Lean-ready interval-cell witnesses, and calibration-safe \(\rho\)/vol-of-vol
projection gates, adds a finite rough-vol path-dependent exotic certificate package with
vanilla-calibrated scenario envelopes, Asian/lookback SOE transfer reserves, barrier shift brackets,
point-model model-risk intervals, Lean-ready rational cell witnesses, and LP dual certificates, adds
a finite rough forward-start smile certificate package with conditional-smile tubes,
martingale-coupling polytopes, LP forward-start and forward-vol envelopes, IV transfer, aggregate
asymptotic coefficient certificates, Lean-ready rational cell witnesses, and rough-model rejection
tests, adds a finite rough-to-classical limit certificate package with Hurst-cell bridges to
\(H=1/2\), kernel/state coupling error budgets, price/IV/Greek continuity bands, Lean-ready rational
cell witnesses, and rough-versus-classical justification gates, adds a finite Lee limsup-to-limit
classification package with declared tail-regularity cells, Lee-slope stabilization bands,
oscillation-witness rejection certificates, finite-strike extrapolation eligibility gates, and
Lean-ready rational witnesses, adds a finite Heston moment-explosion monotonicity package with
right/left explosion-branch interval enclosures, critical \(p^*(T)\)/\(q^*(T)\) bands,
maturity-grid monotonicity signs, convexity diagnostic intervals, Lee-wing sanity gates,
no-false-rejection Heston parameter-cell tests, and Lean-ready rational witnesses, adds a finite
stochastic-time-change tail-transfer package with base-cumulant and clock-domain cells,
martingale-normalization checks, critical-moment transfer brackets, active-boundary labels,
Lee-slope bands, finite-strike slope gates, limit/limsup/rejected/unknown classifications,
calibration-safe no-false-rejection tests, tail-reserve price/IV audit ledgers, and replayable
rational point or Farkas witnesses, adds a finite rough-kernel moment-explosion package with
Volterra--Riccati envelope cells, finite blow-up/non-blow-up brackets, critical
\(p^*(T)\)/\(q^*(T)\) bands, Lee-slope transfer, rough-Heston wing sanity gates,
no-false-rejection rough-affine parameter-cell tests, and replayable rational point or Farkas
witnesses, adds a finite tail-robust calibration objective package with quote/error input packets,
weighted loss ledgers, declared tail-weight cells, no-arbitrage wing-envelope gates,
accepted/rejected tail-cell ledgers, recovery bands for left and right tail parameters,
calibration-to-wing reserve transfer, no-false-rejection desk acceptance rules, and replayable
rational point or Farkas witnesses, adds a finite universal wing-stress package with bounded
wing-budget ledgers, admissible stress directions, LP extrema for wing steepening, flattening, and
book impact, active-facet and dual-multiplier explanations, Lee-slope and variance-reserve handoffs,
accepted/rejected stress ledgers, and replayable rational point, dual, or Farkas witnesses, adds a
finite transport-percentile package with quote/projection input ledgers, quote-to-law extraction
maps, declared ground metrics, maturitywise optimal-transport LP scores, historical regime prototype
distances, empirical tie-rank intervals, DKW confidence bands, repair-sensitivity reserves,
accepted/rejected percentile alerts, and replayable rational point, dual, or Farkas witnesses, adds a
finite joint realized-implied percentile package with paired realized-volatility and IV-surface score
ledgers, declared dependence cells, empirical joint-rank boxes, fixed-cell lower-orthant confidence
bands, measurement and repair reserves, variance-risk-premium alert gates, and replayable rational
point, dual, or Farkas witnesses, adds a finite extreme-percentile stability package with
top/bottom-tail rank buckets, order-statistic threshold bands, tie and score-box reserves,
repair/shock stability margins, abstaining crisis-state alert gates, and replayable rational point,
dual, or Farkas witnesses, adds a finite cross-asset percentile-transfer package with source/target
asset cells, exchange-normalized no-arbitrage surface ledgers, monotone quantile and score-transfer
maps, cross-asset score-box and transport reserves, holdout, DKW, and block-validation confidence
bands, abstaining relative-rich/cheap alert gates, and replayable rational point, dual, or Farkas
witnesses, adds a finite model-error percentile package with calibrated model-family distance
ledgers, per-family and nearest-family residual score boxes, quote/repair/discretization and
calibration reserves, empirical tie-rank intervals, DKW, holdout, and block-validation confidence
bands, abstaining model-break percentile alert gates, and replayable rational point, dual, or
Farkas witnesses, adds a finite causal implied-volatility percentile package with frozen causal
driver cells, adjustment ledgers, counterfactual repaired-surface score boxes, realized-volatility,
jump, rate/dividend, liquidity, and event-shock attribution intervals, DKW, holdout, and
block-validation confidence bands, abstaining no-false-cause alert gates, and replayable rational
point, dual, or Farkas witnesses, adds a finite discrete hedging CLT-certificate package with
declared hedge-error triangular arrays, martingale-difference centering ledgers, conditional
variance and Lindeberg/Berry--Esseen gates, Studentized normal-approximation bands,
transaction-cost, remainder, and jump reserves, abstaining rehedging-risk alert rules, and
replayable rational witnesses, adds a finite jump-corrected gamma-carry package with pre-jump
delta accounting, declared jump/event cells, jump residual and surface-repricing reserves,
event-kernel jump-loss LP bands, after-cost convexity-survival and jump-defeat gates, abstaining
no-false gamma-risk alerts, and replayable rational witnesses, adds a finite portfolio
gamma/vanna/volga aggregation package with common surface-atlas Greeks, joint spot-surface
covariance projection, signed covariance-contracted convexity bands, scenario and moment LP
aggregation, after-cost survival and deficit gates, abstaining no-false book-level alerts, and
replayable rational witnesses, adds finite vanilla-quote gamma-bound certificates with today's
vanilla quote polytope, declared path-class constraints, model-independent gamma P\&L bands,
primal/dual LP attribution, after-cost survival and deficit gates, abstaining no-false robust-gamma
alerts, and replayable rational witnesses, adds a finite gamma-crowding feedback package with
declared dealer-inventory and hedge-flow cells, signed gamma/vanna/volga exposure ledgers,
realized-volatility and surface-drift response intervals, fixed-point/contraction feedback gates,
squeeze/dampening alert rules, abstaining no-false feedback alerts, and replayable rational
witnesses, adds a finite gamma-percentile package with declared historical surface and
realized-volatility state cells, expected gamma-carry score boxes, empirical tie-rank intervals,
DKW, holdout, and block-validation confidence bands, repair, jump, cost, crowding, and selection
reserves, abstaining rich/cheap opportunity alert gates, and replayable rational point, dual, or
Farkas witnesses, adds a finite asynchronous quote-filtering package with timestamped executable
quote ledgers, stale-underlying and spread reserves, bid-ask interval transfer to a common decision
time, quote-feasibility and repair LP/Farkas certificates, handoff gates to the finite IV surface
filter, abstaining no-false intraday surface alerts, and replayable rational point, dual, or Farkas
witnesses, adds a finite microstructure-noise IV percentile package with quote-bounce, spread,
stale-time, repair, and IV-inversion reserve ledgers, covered IV-surface score boxes, empirical
tie-rank intervals, DKW, holdout, and block-validation confidence bands, abstaining no-false
IV-rank alert gates, and replayable rational point, dual, or Farkas witnesses, adds a finite
opening-auction surface-formation equilibrium package with pre-open indication, imbalance,
auction-order, and first-print quote ledgers, finite clearing and no-arbitrage surface cells,
equilibrium residual and selection bands, first-print surface-formation maps, accepted/rejected
formation gates, abstaining no-false auction-surface alerts, and replayable rational point, dual, or
Farkas witnesses, adds a finite adaptive live data-snooping/e-process package with predictable live
strategy menus, nonnegative e-process capital ledgers, optional-stopping and anytime-valid lower
confidence bounds, wealth-split and alpha-spending reuse gates, abstaining no-false live signal
alerts, and replayable rational witnesses, adds a finite adaptive IV-regime segmentation and
conditional percentile package with predictable feature and segment-library ledgers, split/holdout or
earlier-sample segmentation freezes, segment-membership score boxes, empirical tie-rank intervals,
DKW, holdout, and block-validation confidence bands, selection, repair, and feature-membership
reserve ledgers, abstaining rich/cheap conditional-percentile alert gates, and replayable rational
point, dual, or Farkas witnesses, adds a finite dependent IV-percentile confidence-band package with
declared dependence ledgers, score boxes, empirical tie-rank intervals, block, effective-sample,
mixing, or e-process confidence radii, simultaneous confidence rows, abstaining rich/cheap
IV-percentile alert gates, and replayable rational point, dual, or Farkas witnesses, adds a finite
liquidity-adjusted gamma-capacity and position-sizing package with executable-liquidity ledgers,
per-size gamma-carry score boxes, capacity intervals, spread, impact, fill, jump, crowding, and
selection reserves, after-cost survival and deficit gates, abstaining no-false size and capacity
alerts, and replayable rational point, dual, or Farkas witnesses, adds a finite robust
calibration-error allocation and ambiguity-budgeted exotic-pricing package with quote/error budget
ledgers, finite ambiguity balls, exotic price and reserve intervals, multiplier budget attribution,
abstaining no-false acceptance and action gates, and replayable rational point, dual, recession, or
Farkas witnesses, adds a finite quote-outlier/source localization and robust live-calibration
influence package with common-time quote and normalization ledgers, candidate source
counterfactual sets, repaired-surface and calibration-output influence boxes, localization and
rejection gates, abstaining no-false calibration/source alerts, and replayable rational point,
dual, recession, or Farkas witnesses, and
defines the first solution queue for the 500-page options PDF.

\section*{Dependency graph}

Papers 1, 3, 4, and 5 are self-contained (3 and 4 form a pair).  Paper 12's
deterministic discrepancy theorems are logically independent of paper 5, but
paper 12 now cites paper 5 as companion context for the random-banded-Toeplitz
stability obstruction.  Paper 13 is the finite-horizon companion to paper 1 and
cites it.  Within Thread B, paper 14 cites papers 12 (deterministic banded
theorem) and 5 (spectral-edge phenomenon), and paper 16 cites paper 5
(balanced-regime gap).  Within Thread C:
paper 2 is the foundation (proof skeleton, population dichotomy, VC lemma);
papers 6--11 each cite predecessors as ``companion'' manuscripts ---
6 needs 2; 7 needs 2, 6; 8 needs 2, 6, 7; 9 needs 2--8 (citations to 3, 5 of
the thread are to its own numbering: see each paper's bibliography);
10 needs 9 and the skeleton; 11 needs all of 2, 6--10; the paper-15 pair should
be treated as cross-citing companions: both use 11, the general-position paper
cites the curved fixed-family extension, and the curved paper cites the
general-position flat companion.  For submission, Thread
C can be bundled as: (i) paper 2 alone; (ii) papers 6+7+8 as ``critical
scaling, I''; (iii) papers 9+10+11+15 as ``critical scaling, II: stratified
and curved truncations'' --- or kept as is.  Paper 17 is the new Thread D
master file and now carries a repo-local import of the current IV theorem
packets under \texttt{paper17/imported\_results/}; the remaining integration
work is to promote verified imported packets into self-contained theorem
sections and keep candidates with unresolved skeptic gates clearly labeled.

\textbf{19. \texttt{Extremal\_Hard\_Edge\_v1.pdf}} ($\sim$81pp, Thread E).
The extremal hard edge: the minimum least singular value
$M_m=\min_{|T|=m}\sigma_{\min}(A_T)$ over all square row submatrices of a tall
random matrix, generalizing the author's ACHA paper on extreme least singular
values of Gaussian row submatrices and the Balan--Wang phase-retrieval stability
problem (Randomstrasse Open Problem 21).  The base is entropy over small-ball
exponent, $\tfrac1m\log M_m\to-h(\gamma)/d_\F$, and the chapter relaxes each
Gaussian hypothesis by moving one of these two inputs.  \emph{Proved}: an abstract
extremal small-ball (first/second moment) principle; a distribution-free lower
side for essentially every continuous ensemble (bounded or subexponential
density, entrywise density profiles, directional anti-concentration, isotropic
log-concave rows, smoothed deterministic frames, row/column-normalized and
diagonally preconditioned Gaussian frames), with summable-budget almost-sure
lower envelopes and a random-count (Poissonization) transfer; the matching exact
base for all rotationally invariant radially-concentrated ensembles, including a
new exact quaternionic law $M_m=c_\gamma^{-m/4+o_\PP(m)}$ that completes the
division algebras $d_\F\in\{1,2,4\}$; the exact rectangular ($d_\F(r{+}1)$),
$k$-th-singular ($d_\F k^2$), and two-sided (entropy-addition) lower bases; an
upper-tail large-deviation rate; and a deterministic finite-alphabet
projective-pair obstruction whose $\{\pm1\}$ case shows that critical random sign
frames fail real phase-retrieval injectivity with probability one.  The
non-Gaussian \emph{upper} side is reduced --- via a proved two-projection local
central limit theorem (appendix) and a proved high-overlap negligibility estimate
--- to a single overlap moment
$\mathbb E[(1-|\varrho|^2)^{-d_\F/2};\,D\ge\varepsilon m]\le1+o(1)$, which is then
\emph{established} for all iid bounded-density entries by a resolvent-monotonicity
bound (the null projector $uu^*\preceq\eta(GG^*+\eta)^{-1}$) plus an elementary
off-support resolvent estimate --- proved self-contained by a leave-one-out argument
(Sherman--Morrison, the Marchenko--Pastur fixed point, and Hanson--Wright), with the
anisotropic local law (Knowles--Yin; Bloemendal--Erd\H{o}s--Knowles--Yau--Yin) an
optional citation --- with full error control (Efron--Stein variance, Hanson--Wright bias, fixed-point
stability at $-\eta$) for the projected-iid rows of the present setting, so the universality base holds
unconditionally for iid bounded-density entries.  A Gershgorin bound on the
$r$-normal Gram matrix and a Chen--Stein argument then collapse the higher
correlations to two-point terms, reducing the explicit fluctuation law
$1-e^{-e^{d_\F x}}$ (numerically confirmed, Kolmogorov--Smirnov $p\approx0.96$) to a
single critical-scale two-point estimate; the point-process intensity and a
one-sided Gumbel bound are proved outright.  All bases and the Gumbel shape are
confirmed by simulation.  No Toeplitz or structured-symbol content.  Remaining open:
the sharp rate of the self-contained off-support estimate and its extension beyond
bounded density, the critical-scale small-codimension Gumbel residual, the complex phase-retrieval base,
and sharp concentration of the extremal stability constant.  Sources in
\texttt{paper19/}.

\section*{Author verification list (before submission)}

\begin{itemize}
\item Papers 3/4: verify the exact statement of \v{S}tef\'ankov\'a's
inheritance theorem (DCDS 2016) covers the stated form.
\item Paper 2: complete the Berlyand--Jabin--Safsten journal data.
\item Paper 5: verify the Konyagin--Schlag citation page numbers and the
identity quoted from Manrique-Mir\'on's reference [4]; the documented gap
should be communicated to the author of arXiv:2602.12581 before public
posting.
\item Papers 6--11: the companion-paper cross-references currently cite
``manuscript, 2026''; assign arXiv identifiers in dependency order (2, then
6--8, then 9--11).
\item Paper 11: the constants in the witness-box lemma were corrected once
(the $\mathcal Z$-sum bound); the shipped version uses the
$9kh$-budget and is self-consistent.
\item Paper 13: confirm the cited statements of the companion paper~1
(margin-stable depth hierarchy, empty-interior obstruction) and of the
underlying May-2026 manuscript (all-tail semifactor, windowed DC1, metric mean
dimension) match the forms used in the finite-horizon stability lemma and the
coordinate squeeze.
\item Paper 14: the expected-normalized full-bandwidth random model is now a
fixed-sign theorem; keep this clearly separated from the residual
realization-normalized contraction question, and verify the
Bandeira--Boedihardjo--van Handel intrinsic-freeness reference for the
delocalization remark.
\item Paper 15: cross-check the general-position and curved companions against
paper 11 (exponent $2/(2n-j)$); the curved paper states a residual estimate
(curvature-band-uniform constant) as open --- ensure it is labelled, not
claimed.
\item Paper 16: source and 102pp PDF are present in \texttt{paper16/}; before
submission, keep the abstract/index aligned with the proved/conditional split:
finite-section reduction on the zero-winding event, deterministic
clearance-floor uniformization, sharp Gaussian off-endpoint and endpoint
Kac--Rice estimates, the fixed-bulk first-moment probability upper bound, the
sharp global continuum upper constant, and the conditional fixed-bulk/global
first-order sandwiches plus the subcritical negative-moment consequence, the
exact two-endpoint formula and sharp second-order
expansion, the continuum-scale \(O_\tau(N^{-1})\) endpoint contribution, the
continuum-scale \(O_{C_0}(\tau/N)\) natural-endpoint-neighborhood contribution,
the global endpoint floor,
and the sharp uniform
small-ball bound are proved; the
bounded-density fixed-point/grid and speed-conditional local/uniform estimates,
all-moment and subexponential speed-tail verifications, the variance-only
derivative cutoff with quantitative separation/grid detection and
grid-rate bracket, bounded-density
conditioning under circle clearance, and Rademacher endpoint/root-of-unity,
fixed-grid, and bounded-denominator lattice scales are also proved, the
fixed-order \(N^{-1}\)
small-ball scale is shown sharp at a constant lattice radius, and the exact
atom at every fixed primitive \(q\)-th root is shown to have scale
\(\gamma_qN^{-\varphi(q)/2}\) with an explicit kernel-covolume constant exactly
on the admissible parity residue classes, and to vanish on the inadmissible
classes, with the resulting bounded-denominator exact-zero grid floor, the
finite-cyclotomic-layer asymptotic, the primitive quadratic divisibility
classes/constants, the quadratic overlap and second-order corrections, the
degree-four constants, the first two bounded-denominator leading coefficients,
and the full unpeeled bounded-denominator expansion through
\(N^{-2}\), plus the endpoint-inclusive bounded-denominator root-of-unity
expansion through order \(N^{-1}\), also recorded; more generally, the peeled bounded-denominator layer
coefficient is proved after removing all lower cyclotomic degrees. Sharp
constants/two-sided minimum law and continuum non-Gaussian universality at full
bounded-density generality remain open/conditional. For
Gaussian coefficients the balanced conditioning conclusion is unconditional on
the zero-winding event, using Konyagin--Schlag clearance.
\item Paper 17: expand the imported theorem packets and 100-problem atlas into
self-contained theorem sections, keep unresolved candidates labeled, and keep
market applications separate from proved mathematical claims.
\item All papers carry a generative-AI declaration and the author's standing
grant acknowledgments; review both before submission.
\end{itemize}

\section*{Open problems left on the table}

(1) Full conclusions-residuality on the Hilbert cube (paper 1, Problem 7.1);
paper 13 closes the finite-horizon layer and rules out the finite-certificate
method class, but the infinite-horizon question stays open for non-finite-horizon
methods.  (2) Unconditional DC1-genericity on Lipschitz strata (paper 4; rigidity
phenomena documented in the workspace STATE file).  (3) The local limit
theorem and variance asymptotics for $\nu_m(\mathbb D)$ (paper 5).  (4) The
exact critical \emph{constant} for the UDC (paper 7).  (5) The $n=2$ edge case
of paper 9's high-probability statement.  (6) The full-bandwidth
realization-normalized random Toeplitz contraction question --- whether the
worst-case $\sqrt{\log(d/n)}$ is a deterministic artifact after normalizing each
sample to an honest contraction, in the regime \(d/n\to\infty\) (paper 14).
(7) A curvature-band-uniform constant for the
interpolation law near multiply-tangent directions (paper 15, curved
truncations).  (8) Sharpen and universalize paper 16: prove sharp constants and
the two-sided boundary minimum law, prove the polynomial-scale boundary
hypothesis at full bounded-density generality without tail or speed-cutoff
assumptions, and formulate continuum/local-CLT discrete analogues beyond the
fixed-rational, bounded-denominator, and exact-zero grid scales.  (9) Execute
paper 17's options-market program: turn its 100 open problems into solved
theorem packages for arbitrage-free surfaces, IV percentiles, opening prices,
gamma carry, robust pricing, and rough/stochastic volatility.  (10) Paper 19's
keystone, now proved via a resolvent-monotonicity bound plus a self-contained
off-support resolvent estimate (leave-one-out with full error control; the
anisotropic local law optional): obtain the sharp rate $M(D)=1+d_\F/(2D)$ for the overlap moment
$\mathbb E[(1-|\varrho|^2)^{-d_\F/2};\,D\ge\varepsilon m]\le1+o(1)$ of non-Gaussian
isotropic normals, and supply the critical-scale small-codimension two-point estimate
that completes the explicit hard-edge Gumbel fluctuation law; also the complex phase-retrieval exponential base
at the sharp threshold, and sharp concentration of the extremal Balan--Wang
stability constant.

\end{document}

